\documentclass[11pt]{article}

\usepackage{fontspec}
\setmainfont{LibertinusSerif}[Path=D:/FRUITFLY/tools/libertinus/Libertinus-7.051/static/OTF/,Extension=.otf,UprightFont=*-Regular,BoldFont=*-Bold,ItalicFont=*-Italic,BoldItalicFont=*-BoldItalic,SmallCapsFont=*-Regular,SmallCapsFeatures={Letters=SmallCaps}]
\setsansfont{LibertinusSans}[Path=D:/FRUITFLY/tools/libertinus/Libertinus-7.051/static/OTF/,Extension=.otf,UprightFont=*-Regular,BoldFont=*-Bold,ItalicFont=*-Italic]
\setmonofont{LibertinusMono}[Path=D:/FRUITFLY/tools/libertinus/Libertinus-7.051/static/OTF/,Extension=.otf,UprightFont=*-Regular,Scale=0.9]
\usepackage{unicode-math}\setmathfont{LibertinusMath-Regular.otf}[Path=D:/FRUITFLY/tools/libertinus/Libertinus-7.051/static/OTF/]
\usepackage[a4paper,margin=2.3cm]{geometry}
\usepackage[protrusion=true]{microtype}
\usepackage{booktabs,longtable,array,calc,etoolbox,graphicx,xcolor,caption,float}
\usepackage{cite}
\usepackage[hidelinks=false,colorlinks=true,linkcolor=blue!50!black,citecolor=blue!50!black,urlcolor=blue!50!black,pdfauthor={Srihari Muralikrishnan},pdftitle={Supplementary Information},pdfsubject={Computational neuroscience preprint},pdfkeywords={connectome, Drosophila, anoxia, spreading depolarization, reperfusion, bistability, brain death}]{hyperref}
\usepackage{fancyhdr}\pagestyle{fancy}\fancyhf{}\renewcommand{\headrulewidth}{0pt}
\fancyfoot[L]{\footnotesize\color{black!60}Muralikrishnan --- preprint v1.0 (not peer reviewed). Model predictions only.}\fancyfoot[R]{\footnotesize\thepage}
\captionsetup{font=small,labelformat=empty,justification=justified,singlelinecheck=false}
\setlength{\LTcapwidth}{\textwidth}
\newcounter{none}
\AtBeginEnvironment{longtable}{\footnotesize\setlength{\tabcolsep}{3.5pt}}
\setcounter{secnumdepth}{0}
\providecommand{\tightlist}{\setlength{\itemsep}{0pt}\setlength{\parskip}{0pt}}
\makeatletter
\newsavebox\pandoc@box
\newcommand*\pandocbounded[1]{\sbox\pandoc@box{#1}\Gscale@div\@tempa{0.82\textheight}{\dimexpr\ht\pandoc@box+\dp\pandoc@box\relax}\Gscale@div\@tempb{\linewidth}{\wd\pandoc@box}\ifdim\@tempb\p@<\@tempa\p@\let\@tempa\@tempb\fi\ifdim\@tempa\p@<\p@\scalebox{\@tempa}{\usebox\pandoc@box}\else\usebox{\pandoc@box}\fi}
\def\fps@figure{tbp}
\makeatother
\setlength{\parskip}{3pt plus 1pt}\setlength{\parindent}{0pt}
\renewcommand{\refname}{References}
\urlstyle{same}
\graphicspath{{../}}
\emergencystretch=2em
\begin{document}
{\Large\bfseries Supplementary Information\par}\vspace{0.5em}

\textbf{A whole-brain \emph{Drosophila} connectome model predicts a duration- and sign-dependent potassium trap after anoxia and a reset-based route to functional recovery}

Srihari Muralikrishnan (Kalaris Labs). Preprint v1.0, 2026-10-20. Contact: hello@kalarislabs.com

All tables below were extracted programmatically from the analysis result files (\texttt{research/structure/STRUCTURE\_RESULTS.md}, \texttt{research/dynamics/DYNAMICS\_RESULTS.md}, \texttt{research/followup/mechanism/MECHANISM\_RESULTS.md}, \texttt{research/followup/replication/REPLICATION\_RESULTS.md}, \texttt{research/sim/README.md}, \texttt{research/lit/parameters.csv}, \texttt{research/followup/replication/tables/v4\_condition\_summary.csv}, \texttt{research/followup/mechanism/raw/sign\_res/}, \texttt{research/followup/index/tables/}) by \texttt{build\_paper.py}. The claim-by-claim synthesis is \texttt{research/SYNTHESIS.md}. They are not retyped. Every value describes the model, not real tissue.

\subsection{S1. Preregistration (registered 2026-10-16, before any insult simulation)}\label{s1.-preregistration-registered-2026-10-16-before-any-insult-simulation}

The full text is in \texttt{research/lit/preregistration.md}. Summary of hypotheses and falsification criteria:

{\def\LTcaptype{none} % do not increment counter
\begin{longtable}[]{@{}
  >{\raggedright\arraybackslash}p{(\linewidth - 6\tabcolsep) * \real{0.2500}}
  >{\raggedright\arraybackslash}p{(\linewidth - 6\tabcolsep) * \real{0.2500}}
  >{\raggedright\arraybackslash}p{(\linewidth - 6\tabcolsep) * \real{0.2500}}
  >{\raggedright\arraybackslash}p{(\linewidth - 6\tabcolsep) * \real{0.2500}}@{}}
\toprule\noalign{}
\begin{minipage}[b]{\linewidth}\raggedright
H
\end{minipage} & \begin{minipage}[b]{\linewidth}\raggedright
Key observable
\end{minipage} & \begin{minipage}[b]{\linewidth}\raggedright
Supported if
\end{minipage} & \begin{minipage}[b]{\linewidth}\raggedright
Falsified if
\end{minipage} \\
\midrule\noalign{}
\endhead
\bottomrule\noalign{}
\endlastfoot
H1 sharp threshold & FII(f), Reach(f) & logistic preferred by ΔAIC ≤ −10, 10--90\% width ≤ 0.10, irrecoverable beyond f\_c & width \textgreater{} 0.25, or ΔAIC \textgreater{} −2, or recovery at f ≥ f\_c + 0.10 \\
H2 targeted vulnerability & ρ(demand, death time); f\_c,targeted/f\_c,random & ρ ≤ −0.3 and ratio ≤ 0.7 (CI excludes 1) & ρ \textgreater{} −0.1 and ratio CI includes 1 \\
H3 graded reperfusion & ΔFII staged − abrupt & ≥ +0.05 (CI excludes 0) in ≥2/3 severities, null negative control & ΔFII ≤ 0, or benefit survives without excitotoxic term \\
H4 hypothermia & Q10\_window & within {[}0.67, 1.5{]} × model Q10 & D\_max ratio CI includes 1, or outside band \\
H5 identity first & D50 ordering; regime coverage & D50\_FII lowest by ≥0.05; regime ≥5\% of grid & D50\_FII not lowest and coverage \textless{} 1\% \\
H6 replication & outcome concordance & ≥3/4 connectomes match; f\_c max/min ≤ 1.5 & ≤2/4 match, or f\_c max/min \textgreater{} 2 \\
\end{longtable}
}

Operational ``brain death'' (preregistered): \textless1\% of neurons spontaneously active for ≥1 s under baseline drive, no readout response (\textless{} 2 SD of noise), persisting after the protocol ends. As run (deviation S6): spontaneous fraction \textless{} 0.01 and no readout response \textgreater{} 8 Hz after the protocol. This is a model analogue of the clinical construct {[}Greer et al.~2023{]}, not a clinical determination.

\subsection{S2. Simulator constants}\label{s2.-simulator-constants}

\subsubsection{\texorpdfstring{Table S2a. Fast LIF layer and calibration (from \texttt{research/sim/README.md})}{Table S2a. Fast LIF layer and calibration (from research/sim/README.md)}}\label{table-s2a.-fast-lif-layer-and-calibration-from-researchsimreadme.md}

{\def\LTcaptype{none} % do not increment counter
\begin{longtable}[]{@{}
  >{\raggedright\arraybackslash}p{(\linewidth - 4\tabcolsep) * \real{0.3333}}
  >{\raggedright\arraybackslash}p{(\linewidth - 4\tabcolsep) * \real{0.3333}}
  >{\raggedright\arraybackslash}p{(\linewidth - 4\tabcolsep) * \real{0.3333}}@{}}
\toprule\noalign{}
\begin{minipage}[b]{\linewidth}\raggedright
constant
\end{minipage} & \begin{minipage}[b]{\linewidth}\raggedright
value
\end{minipage} & \begin{minipage}[b]{\linewidth}\raggedright
source
\end{minipage} \\
\midrule\noalign{}
\endhead
\bottomrule\noalign{}
\endlastfoot
V\_rest = V\_reset & −52 mV & Shiu et al.~2024 \\
V\_th & −45 mV & Shiu et al.~2024 \\
τ\_m / τ\_syn & 20 ms / 5 ms & Shiu et al.~2024 \\
refractory / delay & 2.2 ms / 1.8 ms & Shiu et al.~2024 \\
w\_syn & 0.275 mV published; \textbf{0.185 mV used} (Princeton ≥5 matrix supercritical at 0.275) & calibrated here \\
Poisson kick & 68.75 mV to v & Shiu et al.~2024 \\
dt & 0.1 ms (fast), 5 ms (slow layer) & --- \\
\end{longtable}
}

{\def\LTcaptype{none} % do not increment counter
\begin{longtable}[]{@{}
  >{\raggedright\arraybackslash}p{(\linewidth - 6\tabcolsep) * \real{0.2500}}
  >{\raggedright\arraybackslash}p{(\linewidth - 6\tabcolsep) * \real{0.2500}}
  >{\raggedright\arraybackslash}p{(\linewidth - 6\tabcolsep) * \real{0.2500}}
  >{\raggedright\arraybackslash}p{(\linewidth - 6\tabcolsep) * \real{0.2500}}@{}}
\toprule\noalign{}
\begin{minipage}[b]{\linewidth}\raggedright
readout (mean Hz)
\end{minipage} & \begin{minipage}[b]{\linewidth}\raggedright
princeton5 w0.275
\end{minipage} & \begin{minipage}[b]{\linewidth}\raggedright
princeton5 w0.185
\end{minipage} & \begin{minipage}[b]{\linewidth}\raggedright
shiu783
\end{minipage} \\
\midrule\noalign{}
\endhead
\bottomrule\noalign{}
\endlastfoot
MN9, sugar 50 / 100 / 200 Hz & 65.6 / 162.4 / 195.6 & 3.3 / 117.7 / 148.3 & 15.8 / 70.6 / 111.8 \\
DNp01 L/R, LC4 100 Hz & 220 / 236 & 213 / 229 & -- \\
DNg100 R, LC9 100 Hz & 124 & 116 & -- \\
DNa02 L/R, left LC9+DNa02 180 Hz & 197 / 0 & n.r. & -- \\
\end{longtable}
}

\subsubsection{\texorpdfstring{Table S2b. Slow metabolic/ionic layer (phenomenological, \textbf{not fitted to \emph{Drosophila} data})}{Table S2b. Slow metabolic/ionic layer (phenomenological, not fitted to Drosophila data)}}\label{table-s2b.-slow-metabolicionic-layer-phenomenological-not-fitted-to-drosophila-data}

{\def\LTcaptype{none} % do not increment counter
\begin{longtable}[]{@{}
  >{\raggedright\arraybackslash}p{(\linewidth - 6\tabcolsep) * \real{0.2500}}
  >{\raggedright\arraybackslash}p{(\linewidth - 6\tabcolsep) * \real{0.2500}}
  >{\raggedright\arraybackslash}p{(\linewidth - 6\tabcolsep) * \real{0.2500}}
  >{\raggedright\arraybackslash}p{(\linewidth - 6\tabcolsep) * \real{0.2500}}@{}}
\toprule\noalign{}
\begin{minipage}[b]{\linewidth}\raggedright
constant
\end{minipage} & \begin{minipage}[b]{\linewidth}\raggedright
default
\end{minipage} & \begin{minipage}[b]{\linewidth}\raggedright
calibrated set (dynamics, mechanism)
\end{minipage} & \begin{minipage}[b]{\linewidth}\raggedright
LHS range (mechanism arm L)
\end{minipage} \\
\midrule\noalign{}
\endhead
\bottomrule\noalign{}
\endlastfoot
τ\_ATP & 20 s & \textbf{30 s} & 15--60 s (log) \\
ATP reserve m & 4 × resting demand & 4 & 2--8 (log) \\
demand exponent β (area\^{}β) & 0.5 & 0.5 & 0.25--1 \\
r\_half (rate demand) & 100 Hz & 100 Hz & --- \\
pump failure e\_ok & 0.75 & 0.75 & --- \\
K₀ / K\_max & 5 / 80 mM & 5 / 80 mM & K\_max 50--100 mM \\
g\_K (Nernst coupling) & 0.6 & 0.6 & 0.3--0.9 \\
k\_rel (SD release) & 60 mM/s & 60 & 20--120 (log) \\
k\_clr (clearance × pump) & 1 s⁻¹ & 1 & 0.5--2 (log) \\
V\_sd & −30 ± 2 mV & same & --- \\
depolarization block & h → 1 − σ((v\_rest,eff + 32)/2), τ\_h 10 ms & same & --- \\
injury hazard & h\_max·J⁴/(J⁴ + J50⁴), h\_max 1/30 s⁻¹ & --- & --- \\
J50 & 60 s & \textbf{600 s} & 300--1,200 s (log) \\
Q10 demand / supply / pump / injury & 2.5 / 1.5 / 1.8 / --- & 2.5 / 1.5 / 1.8 / \textbf{2.5} & not varied \\
speedup (slow-time compression) & --- & \textbf{50} & not varied (25--100 tested in arm A) \\
glutamate sign & inhibitory & inhibitory & inhibitory or excitatory \\
w\_syn & 0.185 & 0.185 & 0.15--0.25 \\
\end{longtable}
}

\subsubsection{Table S2c. Calibration runs (dynamics arm A)}\label{table-s2c.-calibration-runs-dynamics-arm-a}

{\def\LTcaptype{none} % do not increment counter
\begin{longtable}[]{@{}
  >{\raggedright\arraybackslash}p{(\linewidth - 6\tabcolsep) * \real{0.2500}}
  >{\raggedright\arraybackslash}p{(\linewidth - 6\tabcolsep) * \real{0.2500}}
  >{\raggedright\arraybackslash}p{(\linewidth - 6\tabcolsep) * \real{0.2500}}
  >{\raggedright\arraybackslash}p{(\linewidth - 6\tabcolsep) * \real{0.2500}}@{}}
\toprule\noalign{}
\begin{minipage}[b]{\linewidth}\raggedright
condition
\end{minipage} & \begin{minipage}[b]{\linewidth}\raggedright
SD onset (model s)
\end{minipage} & \begin{minipage}[b]{\linewidth}\raggedright
dead
\end{minipage} & \begin{minipage}[b]{\linewidth}\raggedright
FII\_norm
\end{minipage} \\
\midrule\noalign{}
\endhead
\bottomrule\noalign{}
\endlastfoot
tau20\_sp25\_D60\_reset & 9.0 & 0.000 & 1.001 \\
tau20\_sp50\_D60\_reset & 9.5 & 0.000 & 1.001 \\
tau20\_sp100\_D60\_reset & 10.0 & 0.000 & 0.998 \\
tau30\_sp25\_D60\_reset & 11.5 & 0.000 & 1.001 \\
tau30\_sp50\_D60\_reset & 12.5 & 0.000 & 1.001 \\
tau30\_sp100\_D60\_reset & 13.5 & 0.000 & 0.998 \\
tau40\_sp25\_D60\_reset & 14.0 & 0.000 & 1.001 \\
tau40\_sp50\_D60\_reset & 15.5 & 0.000 & 1.001 \\
tau40\_sp100\_D60\_reset & 16.5 & 0.000 & 0.998 \\
J50\_60\_D120\_reset & 12.5 & 0.907 & 0.131 \\
J50\_300\_D120\_reset & 12.5 & 0.017 & 0.951 \\
sham\_naive & -- & 0.000 & 1.001 \\
sham\_reset & -- & 0.000 & 1.001 \\
sp25\_D30\_naive & 11.5 & 0.000 & -0.104 \\
sp100\_D30\_naive & 13.5 & 0.000 & 0.065 \\
\end{longtable}
}

\subsection{\texorpdfstring{S3. Literature parameter table (Table S1; \texttt{research/lit/parameters.csv}, 81 rows, 6 flagged UNVERIFIED)}{S3. Literature parameter table (Table S1; research/lit/parameters.csv, 81 rows, 6 flagged UNVERIFIED)}}\label{s3.-literature-parameter-table-table-s1-researchlitparameters.csv-81-rows-6-flagged-unverified}

Values are as extracted from the cited sources. ``UNVERIFIED'' rows were not used as calibration targets.

{\def\LTcaptype{none} % do not increment counter
\begin{longtable}[]{@{}
  >{\raggedright\arraybackslash}p{(\linewidth - 12\tabcolsep) * \real{0.1429}}
  >{\raggedright\arraybackslash}p{(\linewidth - 12\tabcolsep) * \real{0.1429}}
  >{\raggedright\arraybackslash}p{(\linewidth - 12\tabcolsep) * \real{0.1429}}
  >{\raggedright\arraybackslash}p{(\linewidth - 12\tabcolsep) * \real{0.1429}}
  >{\raggedright\arraybackslash}p{(\linewidth - 12\tabcolsep) * \real{0.1429}}
  >{\raggedright\arraybackslash}p{(\linewidth - 12\tabcolsep) * \real{0.1429}}
  >{\raggedright\arraybackslash}p{(\linewidth - 12\tabcolsep) * \real{0.1429}}@{}}
\toprule\noalign{}
\begin{minipage}[b]{\linewidth}\raggedright
parameter
\end{minipage} & \begin{minipage}[b]{\linewidth}\raggedright
value
\end{minipage} & \begin{minipage}[b]{\linewidth}\raggedright
unit
\end{minipage} & \begin{minipage}[b]{\linewidth}\raggedright
species
\end{minipage} & \begin{minipage}[b]{\linewidth}\raggedright
context
\end{minipage} & \begin{minipage}[b]{\linewidth}\raggedright
source\_doi
\end{minipage} & \begin{minipage}[b]{\linewidth}\raggedright
confidence
\end{minipage} \\
\midrule\noalign{}
\endhead
\bottomrule\noalign{}
\endlastfoot
K\_o\_rest\_fly\_brain & \textasciitilde10 & mM & D. melanogaster (+4 Drosophila spp.) & Resting CNS {[}K+{]}o, K+-selective microelectrode, cold and warm acclimated & 10.1242/jeb.179598 & high (full text) \\
K\_o\_rest\_fly\_brain\_alt & 6.6 ± 0.3 & mM & D. melanogaster w1118 males & Baseline {[}K+{]}o before repetitive N2 anoxia & 10.1242/jeb.074468 & high (abstract) \\
K\_o\_rest\_fly\_brain\_ouabain\_prep & 13.1 ± 2.8 & mM & D. melanogaster & Baseline between ouabain-induced repetitive surges & 10.1371/journal.pone.0028994 & high (full text) \\
K\_o\_SD\_peak\_cold\_warm\_acclimated & 58.3 ± 3.7 & mM & D. melanogaster (25 °C acclimated) & Peak {[}K+{]}o during cold-induced SD & 10.1242/jeb.179598 & high (full text) \\
K\_o\_SD\_peak\_cold\_cold\_acclimated & 44.4 ± 5.7 & mM & D. melanogaster (15 °C acclimated) & Peak {[}K+{]}o during cold-induced SD & 10.1242/jeb.179598 & high (full text) \\
K\_o\_SD\_peak\_species\_range & 29.5 -- 57.9 & mM & Drosophila spp. (montana \ldots{} birchii) & Peak {[}K+{]}o at cold SD; \textasciitilde3--5x resting & 10.1242/jeb.179598 & high (full text) \\
K\_o\_terminal\_plateau & 54.1 ± 13.9 & mM & D. melanogaster & Terminal loss of K+ homeostasis after repetitive ouabain comas (no recovery; death) & 10.1371/journal.pone.0028994 & high (full text) \\
K\_hemolymph\_fold\_change\_2h\_anoxia & \textasciitilde3 & fold & D. melanogaster adult & Hemolymph {[}K+{]} after 2 h N2 anoxia; restored on reoxygenation & 10.1242/jeb.177147 & high (abstract+full text) \\
t\_anoxic\_SD\_onset\_RT & 14.2 ± 0.3 & s & D. melanogaster (25 °C acclimated) & Time to half-amplitude DC drop after pure N2, room temperature & 10.1242/jeb.179598 & high (full text) \\
t\_anoxic\_SD\_onset\_RT\_cold\_acclimated & 16.2 ± 0.6 & s & D. melanogaster (15 °C acclimated) & As above & 10.1242/jeb.179598 & high (full text) \\
t\_anoxic\_SD\_onset\_species\_range & 10.5 -- 19.8 & s & Drosophila spp. & Time to anoxic SD at room temperature & 10.1242/jeb.179598 & high (full text) \\
N2\_pulse\_threshold\_for\_K\_surge & between 5 and 90 (exact value not extracted) & s & D. melanogaster & N2 pulse duration stepped 5 s→90 s revealed a threshold for the {[}K+{]}o surge & 10.1371/journal.pone.0028994 & UNVERIFIED exact value (figure only) \\
t\_paralysis\_anoxia\_adult & \textless30 & s & D. melanogaster adult & Behavioural paralysis after onset of anoxia & 10.1242/jeb.177147 & high (abstract) \\
t\_loss\_of\_coordination\_anoxia & 60 -- 120 & s & D. melanogaster & Loss of coordination/falling in N2 & 10.1016/S0022-1910(96)00084-4 & high (abstract) \\
anoxia\_duration\_with\_full\_recovery & up to 240 & min & D. melanogaster & Flies tolerated N2 up to 4 h and recovered; nonlinear duration--recovery relation & 10.1016/S0022-1910(96)00084-4 & high (abstract) \\
recovery\_time\_after\_\textless2h\_anoxia & \textless=15 & min & D. melanogaster & First movement after reoxygenation (Krishnan 1997 as cited in Campbell 2018) & 10.1242/jeb.177147 & medium (secondary citation) \\
recovery\_time\_after\_4h\_anoxia & \textasciitilde60 & min & D. melanogaster & Mobility recovery time (Krishnan 1997 as cited in Campbell 2018) & 10.1242/jeb.177147 & medium (secondary citation) \\
LT50\_anoxia\_adult & 8.3 & h & D. melanogaster adult & Anoxia survival, 50\% lethal time & 10.1242/jeb.125849 & high (abstract) \\
LT50\_anoxia\_larva & \textasciitilde1 & h & D. melanogaster L3 & Anoxia survival & 10.1242/jeb.125849 & high (abstract) \\
survival\_2h\_anoxia\_adult & 100 & \% & D. melanogaster adult & Despite ATP 3\% of normal & 10.1242/jeb.177147 & high (abstract) \\
ATP\_fraction\_15min\_anoxia & 0.35 & fraction of normoxic & D. melanogaster adult & Whole-body ATP (not brain-specific) & 10.1242/jeb.177147 & high (full text) \\
ATP\_fraction\_30min\_anoxia & \textasciitilde0.03 & fraction of normoxic & D. melanogaster adult & Whole-body ATP; then stable to \textasciitilde8 h, further drop beyond & 10.1242/jeb.177147 & high (full text) \\
tau\_ATP\_depletion\_anoxia & \textasciitilde14 & min & D. melanogaster adult & DERIVED: single-exponential fit to 1.0→0.35 in 15 min (tau=15/ln(1/0.35)); decline is faster than single exponential by 30 min & 10.1242/jeb.177147 & derived (assumption-dependent) \\
metabolic\_rate\_in\_anoxia & \textasciitilde3 & \% of normoxic & D. melanogaster adult and larva & Metabolic suppression when immobile in anoxia & 10.1242/jeb.125849 & high (abstract) \\
O2\_consumption\_low\_O2 & \textasciitilde20 & \% of control & D. melanogaster & O2 consumption in low-O2 atmospheres & 10.1016/S0022-1910(96)00084-4 & high (abstract) \\
T\_cold\_SD\_warm\_acclimated & 5.3 ± 0.5 & °C & D. melanogaster (25 °C acclimated) & Temperature of cold-induced SD, 1 °C/min ramp & 10.1242/jeb.179598 & high (full text) \\
T\_cold\_SD\_cold\_acclimated & 2.4 ± 0.4 & °C & D. melanogaster (15 °C acclimated) & As above & 10.1242/jeb.179598 & high (full text) \\
T\_heat\_coma\_species\_range & 38.69 -- 45.11 & °C & Drosophila spp. (immigrans \ldots{} mojavensis) & Heat coma (Tcoma) under 0.25 °C/min ramp; coincides with CNS SD & 10.1242/jeb.218750 & high (full text, bioRxiv version) \\
T\_heat\_SD\_D\_melanogaster & not extracted & °C & D. melanogaster & Species-specific heat-SD temperature reported in figures/tables only & 10.1242/jeb.218750 & UNVERIFIED \\
static\_heat\_SD\_assay\_temperature & 38 & °C & Drosophila spp. & Constant-temperature exposure producing SD and heat coma within 1 h & 10.1242/jeb.218750 & high (full text) \\
SD\_propagation\_velocity\_RT & 2.5 ± 0.3 & mm/min & D. melanogaster (warm acclimated) & KCl-induced SD, two-electrode measurement & 10.1152/ajpregu.00040.2025 & high (full text) \\
Q10\_SD\_propagation\_velocity & 2.98 (warm acclimated); 2.13 (cold acclimated) & dimensionless & D. melanogaster & Temperature sensitivity of SD propagation & 10.1152/ajpregu.00040.2025 & high (full text) \\
Q10\_SD\_onset\_slope & \textasciitilde1.6 & dimensionless & D. melanogaster & Temperature sensitivity of SD onset slope & 10.1152/ajpregu.00040.2025 & high (full text) \\
Q10\_metabolic\_rate\_general & \textasciitilde2.2 -- 2.4 (at 15--35 °C) & dimensionless & general (ectotherms) & DERIVED from Boltzmann factor with assumed E=0.65 eV; E value not read from the retrieved abstract & 10.1126/science.1061967 & UNVERIFIED (derived) \\
K\_o\_baseline\_drift\_repetitive\_anoxia\_17C & -1.1 ± 1.3 & mM (cumulative Δ) & D. melanogaster w1118 males & Total baseline variation during repetitive anoxia at 17 °C & 10.1242/jeb.074468 & high (abstract) \\
K\_o\_baseline\_drift\_repetitive\_anoxia\_23C & 17.3 ± 1.4 & mM (cumulative Δ) & D. melanogaster w1118 males & As above, 23 °C & 10.1242/jeb.074468 & high (abstract) \\
K\_o\_baseline\_drift\_repetitive\_anoxia\_29C & 332.7 ± 83.0 & mM (cumulative Δ) & D. melanogaster w1118 males & As above, 29 °C (cumulative sum, not a concentration) & 10.1242/jeb.074468 & high (abstract) \\
energy\_share\_action\_potentials & 47 & \% of signalling energy & rodent grey matter & Attwell--Laughlin budget & 10.1097/00004647-200110000-00001 & high (abstract) \\
energy\_share\_postsynaptic & 34 & \% of signalling energy & rodent grey matter & Attwell--Laughlin budget & 10.1097/00004647-200110000-00001 & high (abstract) \\
energy\_share\_resting\_potential & 13 & \% of signalling energy & rodent grey matter & Attwell--Laughlin budget & 10.1097/00004647-200110000-00001 & high (abstract) \\
energy\_share\_cortex\_updated & postsyn 50 / AP 21 / rest 20 / release 5 / recycling 4 & \% of signalling energy & rodent cortex & Updated budget with efficient APs & 10.1038/jcbfm.2012.35 & high (abstract) \\
max\_fraction\_simultaneously\_active & \textless=15 & \% of neurons & rodent cortex (theory) & Energy-limited distributed coding prediction & 10.1097/00004647-200110000-00001 & high (abstract) \\
AP\_energy\_variation\_between\_neuron\_types & \textgreater10 & fold & vertebrate and invertebrate models & Depends on Na+/K+ current overlap & 10.1371/journal.pcbi.1000840 & high (abstract) \\
t\_anoxic\_depolarization\_mammal & 2 -- 5 & min after severe ischaemia onset & mammals (experimental) & Onset of potentially reversible anoxic depolarisation; injury begins with it & 10.1002/ana.25147 & high (abstract) \\
t\_terminal\_SD\_human & 3.9 (IQR 2.6--6.3) & min after final perfusion drop & human (dying patients) & Terminal SD after withdrawal of life support & 10.1002/ana.25147 & high (abstract) \\
ptiO2\_nonspreading\_depression\_human & 11 (IQR 7--14) & mmHg & human & Tissue pO2 at onset of electrical silencing & 10.1002/ana.25147 & high (abstract) \\
CBF\_threshold\_protein\_synthesis & 0.55 & ml/g/min & mammal & Ischaemic flow threshold hierarchy & 10.1002/ana.410360404 & high (abstract) \\
CBF\_threshold\_anaerobic\_glycolysis & 0.35 & ml/g/min & mammal & As above & 10.1002/ana.410360404 & high (abstract) \\
CBF\_threshold\_energy\_failure & \textasciitilde0.20 & ml/g/min & mammal & Transmitter release and energy disturbance begin & 10.1002/ana.410360404 & high (abstract) \\
CBF\_threshold\_anoxic\_depolarization & \textless0.15 & ml/g/min & mammal & Anoxic depolarisation & 10.1002/ana.410360404 & high (abstract) \\
K\_o\_peak\_anoxic\_depolarization\_mammal & \textasciitilde60 (commonly cited) & mM & mammal & Not read in retrieved text (Hansen 1985 closed access) & 10.1152/physrev.1985.65.1.101 & UNVERIFIED \\
window\_irreversible\_injury\_mammal\_normothermic & \textasciitilde5 (classical rule) & min & mammal & Classical estimate; not verified in retrieved text and contested by Hossmann 1973 and BrainEx & 10.1001/archneur.1973.00490300037004 & UNVERIFIED \\
glycolytic\_flux\_increase\_ischaemia & 4 -- 7 & fold & mouse brain & After decapitation ischaemia & 10.1016/S0021-9258(18)51740-3 & high (abstract via record) \\
neuron\_loss\_rate\_stroke & 1.9e6 & neurons/min & human & Untreated large-vessel ischaemic stroke & 10.1161/01.STR.0000196957.55928.ab & high (abstract) \\
stroke\_evolution\_duration & 10 (6--18) & h & human & Onset to infarct completion & 10.1161/01.STR.0000196957.55928.ab & high (abstract) \\
therapeutic\_hypothermia\_target & 32 -- 34 for 24 h & °C & human & After VF cardiac arrest & 10.1056/NEJMoa012689 & high (abstract) \\
hypothermia\_favourable\_outcome & 55 vs 39 (RR 1.40, 95\% CI 1.08--1.81) & \% patients & human & Hypothermia vs normothermia, 6-month neurological outcome & 10.1056/NEJMoa012689 & high (abstract) \\
postconditioning\_effective\_schedule & 3 cycles 30 s reperfusion/10 s occlusion; 10 cycles 10 s/10 s & s & rat (focal ischaemia) & Reduced infarct; gradual reperfusion also protective & 10.1002/jnr.21703 & high (abstract) \\
postconditioning\_onset\_deadline & \textless3 & min after reperfusion & rat & Protection lost when started 3 min after reperfusion & 10.1002/jnr.21703 & high (abstract) \\
postconditioning\_late\_window & up to 6 & h after focal ischaemia & rodent (review) & Reported delayed postconditioning protection & 10.1038/jcbfm.2009.13 & medium (review abstract) \\
BrainEx\_postmortem\_interval & 4 & h & pig (isolated brain) & Warm post-mortem interval before ex vivo perfusion & 10.1038/s41586-019-1099-1 & high (full text) \\
BrainEx\_perfusion\_duration & 6 & h & pig & Normothermic perfusion; cellular/synaptic activity restored, no global ECoG & 10.1038/s41586-019-1099-1 & high (full text) \\
OrganEx\_warm\_ischaemia & 1 & h & pig (whole body) & Warm ischaemia before OrganEx perfusion & 10.1038/s41586-022-05016-1 & high (abstract) \\
LIF\_V\_rest & -52 & mV & D. melanogaster model & Shiu et al.~whole-brain LIF & 10.1038/s41586-024-07763-9 & high (full text) \\
LIF\_V\_reset & -52 & mV & D. melanogaster model & As above & 10.1038/s41586-024-07763-9 & high (full text) \\
LIF\_V\_threshold & -45 & mV & D. melanogaster model & As above & 10.1038/s41586-024-07763-9 & high (full text) \\
LIF\_R\_membrane & 10 & kOhm cm2 & D. melanogaster model & As above & 10.1038/s41586-024-07763-9 & high (full text) \\
LIF\_C\_membrane & 2 & uF/cm2 & D. melanogaster model & As above & 10.1038/s41586-024-07763-9 & high (full text) \\
LIF\_tau\_membrane & 20 & ms & D. melanogaster model & DERIVED: Cm x Rm = 2 uF/cm2 x 10 kOhm cm2 & 10.1038/s41586-024-07763-9 & high (derived from full-text values) \\
LIF\_tau\_synapse & 5 & ms & D. melanogaster model & Synaptic conductance decay & 10.1038/s41586-024-07763-9 & high (full text) \\
LIF\_delay & 1.8 & ms & D. melanogaster model & Spike-to-PSP delay & 10.1038/s41586-024-07763-9 & high (full text) \\
LIF\_refractory & 2.2 & ms & D. melanogaster model & As above & 10.1038/s41586-024-07763-9 & high (full text) \\
LIF\_W\_syn & 0.275 & mV per synapse & D. melanogaster model & Single free parameter & 10.1038/s41586-024-07763-9 & high (full text) \\
LIF\_protocol & 30 trials x 1000 ms; Poisson input to sensory neurons (e.g.~10--200 Hz sugar GRNs); baseline firing 0 Hz & --- & D. melanogaster model & Brian2 implementation & 10.1038/s41586-024-07763-9 & high (full text) \\
FAFB\_neurons & 139255 & neurons & D. melanogaster female brain & FlyWire; \textasciitilde5e7 chemical synapses & 10.1038/s41586-024-07558-y & high (abstract) \\
MCNS\_neurons & 166700 & neurons & D. melanogaster male CNS & 11,710 neuron types & 10.1016/j.cell.2026.08.015 & high (abstract) \\
NaK\_pump\_share\_total\_brain\_ATP & \textasciitilde50 (commonly cited); abstract states only that the great majority of energy goes to signalling, particularly Na+ transport & \% of brain ATP & mammal & Exact share not stated in the retrieved abstract & 10.1016/S0165-0173(00)00038-2 & UNVERIFIED (numeric); qualitative claim verified \\
hypothermia\_Bernard\_protocol & 33 for 12 h & °C & human & Induced within 2 h of ROSC after out-of-hospital cardiac arrest & 10.1056/NEJMoa003289 & high (abstract) \\
hypothermia\_Bernard\_good\_outcome & 49 vs 26 (adj. OR 5.25, 95\% CI 1.47--18.76) & \% patients & human & Good outcome at discharge, hypothermia vs normothermia & 10.1056/NEJMoa003289 & high (abstract) \\
BD\_apnoea\_test\_targets & pH \textless{} 7.30 and PaCO2 \textgreater= 60 & pH / mmHg & human & No spontaneous respiration at these targets (BD/DNC criterion) & 10.1001/jama.2020.11586 & high (abstract) \\
penumbra\_survival\_time & a few & h & rat (focal ischaemia) & Penumbra disappears into core over a few hours & 10.1002/ana.410360404 & medium (qualitative, abstract) \\
\end{longtable}
}

\subsection{S4. Structural percolation (all four connectomes)}\label{s4.-structural-percolation-all-four-connectomes}

\subsubsection{Table S4a. Data and pairing QC}\label{table-s4a.-data-and-pairing-qc}

{\def\LTcaptype{none} % do not increment counter
\begin{longtable}[]{@{}
  >{\raggedright\arraybackslash}p{(\linewidth - 22\tabcolsep) * \real{0.0833}}
  >{\raggedright\arraybackslash}p{(\linewidth - 22\tabcolsep) * \real{0.0833}}
  >{\raggedright\arraybackslash}p{(\linewidth - 22\tabcolsep) * \real{0.0833}}
  >{\raggedright\arraybackslash}p{(\linewidth - 22\tabcolsep) * \real{0.0833}}
  >{\raggedright\arraybackslash}p{(\linewidth - 22\tabcolsep) * \real{0.0833}}
  >{\raggedright\arraybackslash}p{(\linewidth - 22\tabcolsep) * \real{0.0833}}
  >{\raggedright\arraybackslash}p{(\linewidth - 22\tabcolsep) * \real{0.0833}}
  >{\raggedright\arraybackslash}p{(\linewidth - 22\tabcolsep) * \real{0.0833}}
  >{\raggedright\arraybackslash}p{(\linewidth - 22\tabcolsep) * \real{0.0833}}
  >{\raggedright\arraybackslash}p{(\linewidth - 22\tabcolsep) * \real{0.0833}}
  >{\raggedright\arraybackslash}p{(\linewidth - 22\tabcolsep) * \real{0.0833}}
  >{\raggedright\arraybackslash}p{(\linewidth - 22\tabcolsep) * \real{0.0833}}@{}}
\toprule\noalign{}
\begin{minipage}[b]{\linewidth}\raggedright
dataset
\end{minipage} & \begin{minipage}[b]{\linewidth}\raggedright
neurons
\end{minipage} & \begin{minipage}[b]{\linewidth}\raggedright
conn. ids found in annotation
\end{minipage} & \begin{minipage}[b]{\linewidth}\raggedright
annotated ids present in conn.
\end{minipage} & \begin{minipage}[b]{\linewidth}\raggedright
pair-level synapse floor in file
\end{minipage} & \begin{minipage}[b]{\linewidth}\raggedright
pairs (all)
\end{minipage} & \begin{minipage}[b]{\linewidth}\raggedright
pairs ≥5
\end{minipage} & \begin{minipage}[b]{\linewidth}\raggedright
sources (sensory)
\end{minipage} & \begin{minipage}[b]{\linewidth}\raggedright
sinks
\end{minipage} & \begin{minipage}[b]{\linewidth}\raggedright
of which DN / MN
\end{minipage} & \begin{minipage}[b]{\linewidth}\raggedright
removable interneurons
\end{minipage} & \begin{minipage}[b]{\linewidth}\raggedright
modality groups
\end{minipage} \\
\midrule\noalign{}
\endhead
\bottomrule\noalign{}
\endlastfoot
FAFB & 139255 & 100.00\% & 99.52\% & ≥5 & 3732460 & 3732460 & 17550 & 1415 & 1305 / 110 & 120290 & 22 \\
BANC & 158262 & 100.00\% & 97.28\% & ≥3 & 3037361 & 1528585 & 17086 & 2347 & 1316 / 1031 & 138356 & 50 \\
MCNS & 166700 & 100.00\% & 99.48\% & ≥5 & 6242118 & 6242118 & 17937 & 2241 & 1320 / 921 & 146522 & 41 \\
MANC & 23665 & 100.00\% & 99.90\% & ≥1 & 5305638 & 1372588 & 6460 & 838 & 0 / 838 & 16352 & 21 \\
\end{longtable}
}

\subsubsection{Table S4b. Baseline reachability (intact, strict)}\label{table-s4b.-baseline-reachability-intact-strict}

{\def\LTcaptype{none} % do not increment counter
\begin{longtable}[]{@{}
  >{\raggedright\arraybackslash}p{(\linewidth - 12\tabcolsep) * \real{0.1429}}
  >{\raggedright\arraybackslash}p{(\linewidth - 12\tabcolsep) * \real{0.1429}}
  >{\raggedright\arraybackslash}p{(\linewidth - 12\tabcolsep) * \real{0.1429}}
  >{\raggedright\arraybackslash}p{(\linewidth - 12\tabcolsep) * \real{0.1429}}
  >{\raggedright\arraybackslash}p{(\linewidth - 12\tabcolsep) * \real{0.1429}}
  >{\raggedright\arraybackslash}p{(\linewidth - 12\tabcolsep) * \real{0.1429}}
  >{\raggedright\arraybackslash}p{(\linewidth - 12\tabcolsep) * \real{0.1429}}@{}}
\toprule\noalign{}
\begin{minipage}[b]{\linewidth}\raggedright
dataset
\end{minipage} & \begin{minipage}[b]{\linewidth}\raggedright
Reach\_N L4
\end{minipage} & \begin{minipage}[b]{\linewidth}\raggedright
Reach\_N L6
\end{minipage} & \begin{minipage}[b]{\linewidth}\raggedright
Reach\_N ∞
\end{minipage} & \begin{minipage}[b]{\linewidth}\raggedright
Reach\_G L6
\end{minipage} & \begin{minipage}[b]{\linewidth}\raggedright
Reach\_G,loose L6
\end{minipage} & \begin{minipage}[b]{\linewidth}\raggedright
hops to converge
\end{minipage} \\
\midrule\noalign{}
\endhead
\bottomrule\noalign{}
\endlastfoot
FAFB & 0.4132 & 0.8687 & 0.8890 & 0.9808 & 0.9866 & 12 \\
BANC & 0.3940 & 0.7291 & 0.7869 & 0.9417 & 0.9587 & 18 \\
MCNS & 0.6360 & 0.9019 & 0.9394 & 0.9938 & 0.9960 & 10 \\
MANC & 0.9733 & 0.9754 & 0.9754 & 0.9988 & 0.9988 & 10 \\
\end{longtable}
}

\subsubsection{Table S4c. Damage curves, primary metric Reach\_N(L≤6)}\label{table-s4c.-damage-curves-primary-metric-reach_nl6}

{\def\LTcaptype{none} % do not increment counter
\begin{longtable}[]{@{}
  >{\raggedright\arraybackslash}p{(\linewidth - 20\tabcolsep) * \real{0.0909}}
  >{\raggedright\arraybackslash}p{(\linewidth - 20\tabcolsep) * \real{0.0909}}
  >{\raggedright\arraybackslash}p{(\linewidth - 20\tabcolsep) * \real{0.0909}}
  >{\raggedright\arraybackslash}p{(\linewidth - 20\tabcolsep) * \real{0.0909}}
  >{\raggedright\arraybackslash}p{(\linewidth - 20\tabcolsep) * \real{0.0909}}
  >{\raggedright\arraybackslash}p{(\linewidth - 20\tabcolsep) * \real{0.0909}}
  >{\raggedright\arraybackslash}p{(\linewidth - 20\tabcolsep) * \real{0.0909}}
  >{\raggedright\arraybackslash}p{(\linewidth - 20\tabcolsep) * \real{0.0909}}
  >{\raggedright\arraybackslash}p{(\linewidth - 20\tabcolsep) * \real{0.0909}}
  >{\raggedright\arraybackslash}p{(\linewidth - 20\tabcolsep) * \real{0.0909}}
  >{\raggedright\arraybackslash}p{(\linewidth - 20\tabcolsep) * \real{0.0909}}@{}}
\toprule\noalign{}
\begin{minipage}[b]{\linewidth}\raggedright
dataset
\end{minipage} & \begin{minipage}[b]{\linewidth}\raggedright
regime
\end{minipage} & \begin{minipage}[b]{\linewidth}\raggedright
orders
\end{minipage} & \begin{minipage}[b]{\linewidth}\raggedright
f\_c
\end{minipage} & \begin{minipage}[b]{\linewidth}\raggedright
w (10--90\%)
\end{minipage} & \begin{minipage}[b]{\linewidth}\raggedright
ΔAIC
\end{minipage} & \begin{minipage}[b]{\linewidth}\raggedright
best alternative
\end{minipage} & \begin{minipage}[b]{\linewidth}\raggedright
H1 category (prereg)
\end{minipage} & \begin{minipage}[b]{\linewidth}\raggedright
D50
\end{minipage} & \begin{minipage}[b]{\linewidth}\raggedright
f\_c ratio vs random {[}95\% CI{]}
\end{minipage} & \begin{minipage}[b]{\linewidth}\raggedright
D50 ratio vs random {[}95\% CI{]}
\end{minipage} \\
\midrule\noalign{}
\endhead
\bottomrule\noalign{}
\endlastfoot
FAFB & random & 13 & 0.800 & 0.671 & -830.1 & linear & falsified & 0.759 & --- & --- \\
FAFB & metabolic\_area & 10 & 0.277 & 0.421 & -436.2 & exponential & falsified & 0.288 & 0.346 {[}0.337, 0.356{]} & 0.380 {[}0.371, 0.388{]} \\
FAFB & metabolic\_syn & 5 & 0.228 & 0.419 & -253.5 & exponential & falsified & 0.245 & 0.285 {[}0.278, 0.291{]} & 0.323 {[}0.316, 0.330{]} \\
FAFB & degree\_finegrid & 2 & 0.033 & 0.231 & -35.7 & exponential & indeterminate (graded transition) & 0.073 & 0.041 {[}0.041, 0.041{]} & 0.097 {[}0.096, 0.098{]} \\
FAFB & degree & 3 & 0.035 & 0.246 & -40.2 & exponential & indeterminate (graded transition) & 0.075 & 0.044 {[}0.044, 0.044{]} & 0.099 {[}0.098, 0.100{]} \\
FAFB & degree\_matched\_random & 3 & 0.035 & 0.246 & -43.8 & exponential & indeterminate (graded transition) & 0.075 & 0.044 {[}0.044, 0.045{]} & 0.099 {[}0.098, 0.101{]} \\
FAFB & focal & 5 & 0.683 & 1.702 & 3.8 & linear & falsified & 0.690 & 0.854 {[}0.227, 1.008{]} & 0.910 {[}0.662, 1.048{]} \\
BANC & random & 15 & 0.750 & 0.734 & -829.9 & linear & falsified & 0.706 & --- & --- \\
BANC & metabolic\_syn & 8 & 0.107 & 0.232 & -269.3 & exponential & indeterminate (graded transition) & 0.120 & 0.143 {[}0.135, 0.150{]} & 0.170 {[}0.165, 0.176{]} \\
BANC & degree\_finegrid & 2 & 0.017 & 0.086 & -84.7 & exponential & supported (structural criteria) & 0.028 & 0.022 {[}0.022, 0.023{]} & 0.040 {[}0.040, 0.041{]} \\
BANC & degree & 3 & 0.017 & 0.083 & -147.0 & exponential & supported (structural criteria) & 0.030 & 0.023 {[}0.022, 0.024{]} & 0.042 {[}0.041, 0.043{]} \\
BANC & degree\_matched\_random & 2 & 0.019 & 0.081 & -120.4 & exponential & supported (structural criteria) & 0.031 & 0.025 {[}0.024, 0.025{]} & 0.043 {[}0.043, 0.044{]} \\
BANC & focal & 4 & 0.031 & 0.337 & 2.3 & exponential & falsified & 0.283 & 0.041 {[}0.038, 0.913{]} & 0.401 {[}0.134, 0.912{]} \\
MCNS & random & 6 & 0.951 (extrapolated) & 0.744 & -184.4 & linear & falsified & 0.869 & --- & --- \\
MCNS & metabolic\_syn & 6 & 0.304 & 0.440 & -242.4 & exponential & falsified & 0.316 & 0.319 {[}0.309, 0.329{]} & 0.364 {[}0.358, 0.368{]} \\
MCNS & degree\_finegrid & 2 & 0.058 & 0.378 & -2.0 & exponential & falsified & 0.115 & 0.061 {[}0.059, 0.062{]} & 0.133 {[}0.131, 0.134{]} \\
MCNS & degree & 2 & 0.076 & 0.306 & -40.1 & exponential & falsified & 0.115 & 0.080 {[}0.077, 0.083{]} & 0.132 {[}0.131, 0.134{]} \\
MCNS & degree\_matched\_random & 2 & 0.072 & 0.327 & -41.1 & exponential & falsified & 0.118 & 0.076 {[}0.074, 0.078{]} & 0.135 {[}0.134, 0.136{]} \\
MCNS & focal & 3 & 0.028 & 0.521 & 2.6 & exponential & falsified & 0.368 & 0.029 {[}0.029, 0.997{]} & 0.423 {[}0.137, 0.990{]} \\
MANC & random & 20 & 1.027 (extrapolated) & 0.393 & -1045.2 & linear & falsified & \textgreater0.9 & --- & --- \\
MANC & metabolic\_syn & 20 & 0.677 & 0.315 & -1805.5 & linear & falsified & 0.673 & 0.660 {[}0.644, 0.672{]} & \textless{} 0.748 (random never fell to 50\%) \\
MANC & degree\_finegrid & 2 & 0.421 (extrapolated beyond fine-grid max f = 0.3) & 0.467 & -142.4 & linear & falsified & \textgreater0.3 & 0.410 {[}0.402, 0.418{]} & n/a (both censored) \\
MANC & degree & 5 & 0.388 & 0.412 & -572.5 & exponential & falsified & 0.384 & 0.377 {[}0.369, 0.384{]} & \textless{} 0.426 (random never fell to 50\%) \\
MANC & degree\_matched\_random & 5 & 0.394 & 0.421 & -632.1 & linear & falsified & 0.396 & 0.384 {[}0.375, 0.391{]} & \textless{} 0.440 (random never fell to 50\%) \\
MANC & focal & 10 & 0.847 & 0.660 & -86.9 & linear & falsified & 0.727 & 0.825 {[}0.758, 0.917{]} & \textless{} 0.808 (random never fell to 50\%) \\
\end{longtable}
}

\subsubsection{Table S4d. D50 (f at 50\% of baseline) per structural metric}\label{table-s4d.-d50-f-at-50-of-baseline-per-structural-metric}

{\def\LTcaptype{none} % do not increment counter
\begin{longtable}[]{@{}
  >{\raggedright\arraybackslash}p{(\linewidth - 16\tabcolsep) * \real{0.1111}}
  >{\raggedright\arraybackslash}p{(\linewidth - 16\tabcolsep) * \real{0.1111}}
  >{\raggedright\arraybackslash}p{(\linewidth - 16\tabcolsep) * \real{0.1111}}
  >{\raggedright\arraybackslash}p{(\linewidth - 16\tabcolsep) * \real{0.1111}}
  >{\raggedright\arraybackslash}p{(\linewidth - 16\tabcolsep) * \real{0.1111}}
  >{\raggedright\arraybackslash}p{(\linewidth - 16\tabcolsep) * \real{0.1111}}
  >{\raggedright\arraybackslash}p{(\linewidth - 16\tabcolsep) * \real{0.1111}}
  >{\raggedright\arraybackslash}p{(\linewidth - 16\tabcolsep) * \real{0.1111}}
  >{\raggedright\arraybackslash}p{(\linewidth - 16\tabcolsep) * \real{0.1111}}@{}}
\toprule\noalign{}
\begin{minipage}[b]{\linewidth}\raggedright
dataset
\end{minipage} & \begin{minipage}[b]{\linewidth}\raggedright
regime
\end{minipage} & \begin{minipage}[b]{\linewidth}\raggedright
weighted influence
\end{minipage} & \begin{minipage}[b]{\linewidth}\raggedright
GSCC
\end{minipage} & \begin{minipage}[b]{\linewidth}\raggedright
Reach\_N L4
\end{minipage} & \begin{minipage}[b]{\linewidth}\raggedright
Reach\_N L6
\end{minipage} & \begin{minipage}[b]{\linewidth}\raggedright
Reach\_N ∞
\end{minipage} & \begin{minipage}[b]{\linewidth}\raggedright
redundancy
\end{minipage} & \begin{minipage}[b]{\linewidth}\raggedright
Reach\_G L6
\end{minipage} \\
\midrule\noalign{}
\endhead
\bottomrule\noalign{}
\endlastfoot
FAFB & random & 0.497 & 0.506 & 0.626 & 0.759 & 0.853 & 0.858 & \textgreater0.9 \\
FAFB & metabolic\_syn & 0.162 & 0.414 & 0.133 & 0.245 & 0.437 & 0.302 & 0.436 \\
FAFB & metabolic\_area & 0.188 & 0.433 & 0.167 & 0.288 & 0.468 & 0.347 & 0.450 \\
FAFB & degree\_finegrid & 0.086 & \textgreater0.3 & 0.031 & 0.073 & 0.266 & 0.118 & 0.195 \\
BANC & random & 0.295 & 0.417 & 0.554 & 0.706 & 0.825 & 0.804 & 0.862 \\
BANC & metabolic\_syn & 0.065 & 0.183 & 0.071 & 0.120 & 0.218 & 0.173 & 0.230 \\
BANC & degree\_finegrid & 0.031 & 0.108 & 0.014 & 0.028 & 0.065 & 0.055 & 0.074 \\
MCNS & random & 0.327 & 0.528 & 0.806 & 0.869 (1 cens.) & \textgreater0.9 & 0.896 (1 cens.) & \textgreater0.9 \\
MCNS & metabolic\_syn & 0.091 & 0.462 & 0.204 & 0.316 & 0.503 & 0.357 & 0.490 \\
MCNS & degree\_finegrid & 0.057 & \textgreater0.3 & 0.056 & 0.115 & 0.275 & 0.175 & 0.237 \\
MANC & random & 0.240 & 0.656 & 0.884 (1 cens.) & \textgreater0.9 & \textgreater0.9 & \textgreater0.9 & \textgreater0.9 \\
MANC & metabolic\_syn & 0.133 & 0.584 & 0.516 & 0.673 & 0.759 & 0.634 & 0.824 \\
MANC & degree\_finegrid & 0.091 & \textgreater0.3 & 0.223 & \textgreater0.3 & \textgreater0.3 & \textgreater0.3 & \textgreater0.3 \\
\end{longtable}
}

\textbf{Controls.} Degree-preserving rewired null: FAFB random f\_c 0.893 (real 0.800), hub-targeted 0.266 (real 0.035); MANC random 1.107 (extrapolated; real 1.027), hub-targeted 0.613 (real 0.388). Edge cut-off ≥10: FAFB random f\_c 0.607, w 0.960; hub-targeted f\_c 0.013, w 0.087, ΔAIC 0.2 (falsified). MANC ≥1: random f\_c 0.787, w 0.124 (indeterminate); hub-targeted f\_c 0.770. The H1 category for the hub-targeted regime (FAFB) and the random regime (MANC) is therefore not robust to the edge cut-off. Negative control (off-path neurons removed): Δ = 0 for every metric in every dataset (off-path fractions FAFB 2.1\%, BANC 27.0\%, MCNS 2.4\%, MANC 3.2\%).

\textbf{Keystones.} Max relative ΔReach\_N(L6) for single-type removal: FAFB 0.88\% (CB1231; 154 types evaluated), BANC 1.03\% (IN13B010; 218), MCNS 0.86\% (AN04A001; 194), MANC 0.00\% (161). Top-50 membership overlap did not exceed chance once the universe was restricted to types evaluated in both datasets (p = 0.129--0.634). Graded importance replicated (Spearman ρ of influence loss: FAFB--BANC 0.730, FAFB--MCNS 0.710, BANC--MCNS 0.911, BANC--MANC 0.652, MCNS--MANC 0.776).

\begin{figure}
\centering
\includegraphics[width=1\linewidth,height=\textheight,keepaspectratio,alt={Figure S1. (a) FAFB random loss, all structural metrics; (b, c) degree-preserving null and edge cut-off controls for FAFB and MANC.}]{figures/FigS1_structure_metrics_controls.png}
\caption{\textbf{Figure S1.} (a) FAFB random loss, all structural metrics; (b, c) degree-preserving null and edge cut-off controls for FAFB and MANC.}
\end{figure}

\begin{figure}
\centering
\includegraphics[width=1\linewidth,height=\textheight,keepaspectratio,alt={Figure S2. (a) Relative reach lost when all neurons of one cell type are removed (each dot a type; bar = median); (b--d) cross-dataset agreement of relative weighted-influence loss for shared types (Spearman ρ).}]{figures/FigS2_keystone_types.png}
\caption{\textbf{Figure S2.} (a) Relative reach lost when all neurons of one cell type are removed (each dot a type; bar = median); (b--d) cross-dataset agreement of relative weighted-influence loss for shared types (Spearman ρ).}
\end{figure}

\subsection{S5. Dynamics: brain-death map, protocols, tolerance, neuroprotection}\label{s5.-dynamics-brain-death-map-protocols-tolerance-neuroprotection}

\subsubsection{Table S5a. Recoverability boundary (20-s block reset)}\label{table-s5a.-recoverability-boundary-20-s-block-reset}

{\def\LTcaptype{none} % do not increment counter
\begin{longtable}[]{@{}
  >{\raggedright\arraybackslash}p{(\linewidth - 10\tabcolsep) * \real{0.1667}}
  >{\raggedright\arraybackslash}p{(\linewidth - 10\tabcolsep) * \real{0.1667}}
  >{\raggedright\arraybackslash}p{(\linewidth - 10\tabcolsep) * \real{0.1667}}
  >{\raggedright\arraybackslash}p{(\linewidth - 10\tabcolsep) * \real{0.1667}}
  >{\raggedright\arraybackslash}p{(\linewidth - 10\tabcolsep) * \real{0.1667}}
  >{\raggedright\arraybackslash}p{(\linewidth - 10\tabcolsep) * \real{0.1667}}@{}}
\toprule\noalign{}
\begin{minipage}[b]{\linewidth}\raggedright
T (°C)
\end{minipage} & \begin{minipage}[b]{\linewidth}\raggedright
SD onset (s)
\end{minipage} & \begin{minipage}[b]{\linewidth}\raggedright
D\_max (model s)
\end{minipage} & \begin{minipage}[b]{\linewidth}\raggedright
dead at 240 s
\end{minipage} & \begin{minipage}[b]{\linewidth}\raggedright
outcome at 480 s
\end{minipage} & \begin{minipage}[b]{\linewidth}\raggedright
brain death from
\end{minipage} \\
\midrule\noalign{}
\endhead
\bottomrule\noalign{}
\endlastfoot
17 & 15.0 & 359 {[}342, 381{]} & 0.5 \% & structural loss (16.8 \% dead, FII\_norm 0.68) & not reached (960 s: 98.7 \% dead, survivors in runaway) \\
22 & 12.5 & 271 {[}261, 278{]} & 3.8 \% & structural loss (66 \% dead) & 960 s (100 \% dead) \\
29 & 8.5 & 149 {[}145, 153{]} & 40 \% & brain death (99.7 \% dead) & 480 s \\
33 & 6.5 & 124 {[}112, 129{]} & 81 \% & brain death & 480 s \\
\end{longtable}
}

\subsubsection{Table S5b. All 32 reperfusion protocols, ranked by mean FII\_norm (30/180/360 s anoxia, 5 seeds)}\label{table-s5b.-all-32-reperfusion-protocols-ranked-by-mean-fii_norm-30180360-s-anoxia-5-seeds}

{\def\LTcaptype{none} % do not increment counter
\begin{longtable}[]{@{}
  >{\raggedright\arraybackslash}p{(\linewidth - 12\tabcolsep) * \real{0.1429}}
  >{\raggedright\arraybackslash}p{(\linewidth - 12\tabcolsep) * \real{0.1429}}
  >{\raggedright\arraybackslash}p{(\linewidth - 12\tabcolsep) * \real{0.1429}}
  >{\raggedright\arraybackslash}p{(\linewidth - 12\tabcolsep) * \real{0.1429}}
  >{\raggedright\arraybackslash}p{(\linewidth - 12\tabcolsep) * \real{0.1429}}
  >{\raggedright\arraybackslash}p{(\linewidth - 12\tabcolsep) * \real{0.1429}}
  >{\raggedright\arraybackslash}p{(\linewidth - 12\tabcolsep) * \real{0.1429}}@{}}
\toprule\noalign{}
\begin{minipage}[b]{\linewidth}\raggedright
rank
\end{minipage} & \begin{minipage}[b]{\linewidth}\raggedright
protocol
\end{minipage} & \begin{minipage}[b]{\linewidth}\raggedright
30 s
\end{minipage} & \begin{minipage}[b]{\linewidth}\raggedright
180 s
\end{minipage} & \begin{minipage}[b]{\linewidth}\raggedright
360 s
\end{minipage} & \begin{minipage}[b]{\linewidth}\raggedright
dead @360 s
\end{minipage} & \begin{minipage}[b]{\linewidth}\raggedright
outcome 30/180/360
\end{minipage} \\
\midrule\noalign{}
\endhead
\bottomrule\noalign{}
\endlastfoot
1 & revive\_reset & 1.001 & 1.001 & 1.001 & 0.000 & recovered/recovered/recovered \\
2 & clr3\_reset & 1.001 & 0.963 & 0.588 & 0.244 & recovered/recovered/structural\_loss \\
3 & noise\_reset & 1.001 & 0.960 & 0.585 & 0.236 & recovered/recovered/structural\_loss \\
4 & adapt\_reset & 1.001 & 0.959 & 0.577 & 0.248 & recovered/recovered/structural\_loss \\
5 & block20\_del5 & 1.001 & 0.959 & 0.577 & 0.248 & recovered/recovered/structural\_loss \\
6 & staged\_fwd & 1.001 & 0.959 & 0.577 & 0.248 & recovered/recovered/structural\_loss \\
7 & block40 & 1.001 & 0.959 & 0.577 & 0.248 & recovered/recovered/structural\_loss \\
8 & block20 & 1.001 & 0.953 & 0.578 & 0.248 & recovered/recovered/structural\_loss \\
9 & staged\_rev & 1.001 & 0.953 & 0.578 & 0.248 & recovered/recovered/structural\_loss \\
10 & sdblock\_reset & 1.001 & 0.953 & 0.577 & 0.248 & recovered/recovered/structural\_loss \\
11 & block20\_del20 & 1.001 & 0.957 & 0.566 & 0.250 & recovered/recovered/structural\_loss \\
12 & block10 & 1.001 & 0.959 & 0.428 & 0.248 & recovered/recovered/runaway \\
13 & adapt & 0.734 & 0.731 & 0.548 & 0.248 & runaway/runaway/runaway \\
14 & voff-10 & 0.155 & 0.361 & 0.464 & 0.246 & runaway/runaway/runaway \\
15 & voff-20 & 0.137 & 0.320 & 0.435 & 0.243 & runaway/runaway/runaway \\
16 & sdblock & 0.409 & 0.010 & 0.319 & 0.252 & runaway/runaway/runaway \\
17 & clr3 & 0.718 & -0.170 & 0.020 & 0.244 & runaway/runaway/runaway \\
18 & ramp10 & 0.402 & 0.161 & -0.157 & 0.267 & runaway/runaway/runaway \\
19 & gaba4 & 0.183 & 0.349 & -0.173 & 0.239 & runaway/runaway/runaway \\
20 & noise & 0.260 & 0.015 & 0.008 & 0.240 & runaway/runaway/runaway \\
21 & naive & -0.058 & 0.084 & 0.085 & 0.252 & runaway/runaway/runaway \\
22 & gaba2 & 0.191 & 0.100 & -0.195 & 0.239 & runaway/runaway/runaway \\
23 & block5 & -0.060 & 0.097 & 0.018 & 0.252 & runaway/runaway/runaway \\
24 & block2 & -0.088 & 0.098 & 0.010 & 0.252 & runaway/runaway/runaway \\
25 & naive\_noinj & -0.079 & -0.078 & 0.105 & 0.250 & runaway/runaway/runaway \\
26 & pulsed6x5 & -0.104 & -0.032 & -0.123 & 0.281 & runaway/runaway/runaway \\
27 & pulsed3x10 & -0.057 & -0.093 & -0.136 & 0.273 & runaway/runaway/runaway \\
28 & ramp60 & -0.104 & -0.065 & -0.212 & 0.318 & runaway/runaway/runaway \\
29 & ramp30\_noinj & -0.305 & 0.064 & -0.175 & 0.281 & runaway/runaway/runaway \\
30 & ramp30 & -0.301 & 0.063 & -0.246 & 0.292 & runaway/runaway/runaway \\
31 & cool12 & -0.161 & -0.127 & -0.227 & 0.260 & runaway/runaway/runaway \\
32 & cool17 & -0.229 & -0.305 & -0.161 & 0.256 & runaway/runaway/runaway \\
\end{longtable}
}

Protocol key: \texttt{naive} = abrupt step to S = 1; \texttt{blockN} = total spike + synaptic block for N s from reperfusion; \texttt{block20\_delN} = block starting N s after reperfusion; \texttt{rampN} = linear supply ramp over N s; \texttt{pulsedAxB} = A cycles of B-s reperfusion; \texttt{voff-N} = −N mV offset; \texttt{gabaN} = N× GABAergic gain; \texttt{sdblock} = k\_rel × 0.2; \texttt{clr3} = 3× clearance; \texttt{noise} = low-level background Poisson; \texttt{adapt} = spike-frequency adaptation k\_adapt 0.2; \texttt{coolN} = cooling to N °C during reperfusion; \texttt{staged\_fwd/rev} = staged release sensory → central → descending (or reverse), 10 s per stage after a 20-s block; \texttt{revive\_reset} = counterfactual revival of dead neurons + reset; \texttt{\_noinj} = injury off; \texttt{\_reset} suffix = followed by 20-s block.

\subsubsection{Table S5c. Tolerance and preconditioning (arm D; reset protocol)}\label{table-s5c.-tolerance-and-preconditioning-arm-d-reset-protocol}

{\def\LTcaptype{none} % do not increment counter
\begin{longtable}[]{@{}
  >{\raggedright\arraybackslash}p{(\linewidth - 10\tabcolsep) * \real{0.1667}}
  >{\raggedright\arraybackslash}p{(\linewidth - 10\tabcolsep) * \real{0.1667}}
  >{\raggedright\arraybackslash}p{(\linewidth - 10\tabcolsep) * \real{0.1667}}
  >{\raggedright\arraybackslash}p{(\linewidth - 10\tabcolsep) * \real{0.1667}}
  >{\raggedright\arraybackslash}p{(\linewidth - 10\tabcolsep) * \real{0.1667}}
  >{\raggedright\arraybackslash}p{(\linewidth - 10\tabcolsep) * \real{0.1667}}@{}}
\toprule\noalign{}
\begin{minipage}[b]{\linewidth}\raggedright
condition
\end{minipage} & \begin{minipage}[b]{\linewidth}\raggedright
SD onset (s)
\end{minipage} & \begin{minipage}[b]{\linewidth}\raggedright
dead 180 s
\end{minipage} & \begin{minipage}[b]{\linewidth}\raggedright
FII\_norm 180 s
\end{minipage} & \begin{minipage}[b]{\linewidth}\raggedright
dead 360 s
\end{minipage} & \begin{minipage}[b]{\linewidth}\raggedright
FII\_norm 360 s
\end{minipage} \\
\midrule\noalign{}
\endhead
\bottomrule\noalign{}
\endlastfoot
hypomet0.2 & 33.0 & 0.001 & 0.995 & 0.088 & 0.704 {[}0.573, 0.784{]} \\
hypomet0.5 & 19.5 & 0.005 & 0.983 & 0.198 & 0.631 {[}0.576, 0.662{]} \\
hypomet0.8 & 14.5 & 0.008 & 0.970 & 0.235 & 0.568 {[}0.452, 0.657{]} \\
precool12 & 17.0 & 0.000 & 1.001 & 0.006 & 0.977 {[}0.966, 0.997{]} \\
precool17 & 15.0 & 0.001 & 0.985 & 0.041 & 0.908 {[}0.890, 0.938{]} \\
precond15 & 12.5 & 0.009 & 0.961 & 0.249 & 0.591 {[}0.442, 0.681{]} \\
arrest0.05 & 12.5 & 0.001 & 0.983 & 0.059 & 0.711 {[}0.608, 0.785{]} \\
unprotected (C block20) & 12.5 & 0.009 & 0.953 & 0.248 & 0.578 \\
\end{longtable}
}

\subsubsection{Table S5d. Targeted (hub) neuroprotection (arm H; 360-s anoxia)}\label{table-s5d.-targeted-hub-neuroprotection-arm-h-360-s-anoxia}

{\def\LTcaptype{none} % do not increment counter
\begin{longtable}[]{@{}
  >{\raggedright\arraybackslash}p{(\linewidth - 10\tabcolsep) * \real{0.1667}}
  >{\raggedright\arraybackslash}p{(\linewidth - 10\tabcolsep) * \real{0.1667}}
  >{\raggedright\arraybackslash}p{(\linewidth - 10\tabcolsep) * \real{0.1667}}
  >{\raggedright\arraybackslash}p{(\linewidth - 10\tabcolsep) * \real{0.1667}}
  >{\raggedright\arraybackslash}p{(\linewidth - 10\tabcolsep) * \real{0.1667}}
  >{\raggedright\arraybackslash}p{(\linewidth - 10\tabcolsep) * \real{0.1667}}@{}}
\toprule\noalign{}
\begin{minipage}[b]{\linewidth}\raggedright
protected set
\end{minipage} & \begin{minipage}[b]{\linewidth}\raggedright
n neurons
\end{minipage} & \begin{minipage}[b]{\linewidth}\raggedright
FII\_norm (reset)
\end{minipage} & \begin{minipage}[b]{\linewidth}\raggedright
dead
\end{minipage} & \begin{minipage}[b]{\linewidth}\raggedright
FII gain per 1,000 protected
\end{minipage} & \begin{minipage}[b]{\linewidth}\raggedright
FII\_norm (naive)
\end{minipage} \\
\midrule\noalign{}
\endhead
\bottomrule\noalign{}
\endlastfoot
top-degree 1\% & 1393 & 0.703 {[}0.668, 0.743{]} & 0.245 & 0.090 & -0.218 \\
random 1\% & 1393 & 0.568 {[}0.530, 0.591{]} & 0.246 & -0.007 & 0.236 \\
top-degree 2\% & 2785 & 0.744 {[}0.714, 0.791{]} & 0.243 & 0.059 & -0.254 \\
random 2\% & 2785 & 0.590 {[}0.559, 0.617{]} & 0.243 & 0.004 & 0.172 \\
top-degree 5\% & 6963 & 0.847 {[}0.789, 0.905{]} & 0.235 & 0.039 & -0.253 \\
random 5\% & 6963 & 0.574 {[}0.531, 0.596{]} & 0.236 & -0.001 & 0.117 \\
top-degree 10\% & 13926 & 0.906 {[}0.845, 0.944{]} & 0.222 & 0.024 & 0.342 \\
random 10\% & 13926 & 0.588 {[}0.563, 0.637{]} & 0.224 & 0.001 & -0.118 \\
whole brain & 139255 & 1.001 {[}1.001, 1.001{]} & 0.000 & 0.0030 & -0.041 \\
none (C block20 / naive) & 0 & 0.578 & 0.248 & -- & 0.085 \\
\end{longtable}
}

\subsection{S6. Mechanism follow-up}\label{s6.-mechanism-follow-up}

\subsubsection{Table S6a. Latin-hypercube ranges (48 sets + calibrated set)}\label{table-s6a.-latin-hypercube-ranges-48-sets-calibrated-set}

{\def\LTcaptype{none} % do not increment counter
\begin{longtable}[]{@{}ll@{}}
\toprule\noalign{}
parameter & range \\
\midrule\noalign{}
\endhead
\bottomrule\noalign{}
\endlastfoot
w\_syn & 0.15--0.25 \\
τ\_ATP & 15--60 s, log scale \\
J50 & 300--1,200 s, log scale \\
g\_K & 0.3--0.9 \\
k\_rel & 20--120, log scale \\
k\_clr & 0.5--2, log scale \\
reserve & 2--8, log scale \\
demand β & 0.25--1 \\
K\_max & 50--100 mM \\
glutamate & inhibitory or excitatory \\
\end{longtable}
}

\subsubsection{Table S6b. Runaway and rescue across parameter sets}\label{table-s6b.-runaway-and-rescue-across-parameter-sets}

{\def\LTcaptype{none} % do not increment counter
\begin{longtable}[]{@{}
  >{\raggedright\arraybackslash}p{(\linewidth - 8\tabcolsep) * \real{0.2000}}
  >{\raggedright\arraybackslash}p{(\linewidth - 8\tabcolsep) * \real{0.2000}}
  >{\raggedright\arraybackslash}p{(\linewidth - 8\tabcolsep) * \real{0.2000}}
  >{\raggedright\arraybackslash}p{(\linewidth - 8\tabcolsep) * \real{0.2000}}
  >{\raggedright\arraybackslash}p{(\linewidth - 8\tabcolsep) * \real{0.2000}}@{}}
\toprule\noalign{}
\begin{minipage}[b]{\linewidth}\raggedright
anoxia
\end{minipage} & \begin{minipage}[b]{\linewidth}\raggedright
protocol
\end{minipage} & \begin{minipage}[b]{\linewidth}\raggedright
sets with runaway (majority of seeds)
\end{minipage} & \begin{minipage}[b]{\linewidth}\raggedright
sets rescued (FII\_norm ≥ 0.8)
\end{minipage} & \begin{minipage}[b]{\linewidth}\raggedright
mean FII\_norm {[}95 \% CI{]}
\end{minipage} \\
\midrule\noalign{}
\endhead
\bottomrule\noalign{}
\endlastfoot
30 s & abrupt & \textbf{48/48 = 1.00 {[}1.00, 1.00{]}} & 4/48 = 0.08 {[}0.02, 0.17{]} & 0.24 {[}0.14, 0.34{]} \\
30 s & abrupt + 20-s block & 0/48 & \textbf{48/48 = 1.00 {[}1.00, 1.00{]}} & 1.001 {[}1.000, 1.001{]} \\
180 s & abrupt & \textbf{48/48 = 1.00 {[}1.00, 1.00{]}} & 2/48 = 0.04 {[}0.00, 0.10{]} & 0.17 {[}0.08, 0.27{]} \\
180 s & abrupt + 20-s block & 11/48 = 0.23 {[}0.13, 0.35{]} & \textbf{34/48 = 0.71 {[}0.58, 0.83{]}} & 0.75 {[}0.63, 0.85{]} \\
\end{longtable}
}

\subsubsection{Table S6c. Excitotoxic term calibration (180-s anoxia)}\label{table-s6c.-excitotoxic-term-calibration-180-s-anoxia}

{\def\LTcaptype{none} % do not increment counter
\begin{longtable}[]{@{}lll@{}}
\toprule\noalign{}
k\_exc & dead, abrupt & dead, + block \\
\midrule\noalign{}
\endhead
\bottomrule\noalign{}
\endlastfoot
0 & 0.9 \% & 0.9 \% \\
0.3 & 0.9 \% & 0.9 \% \\
1 & 1.1 \% & 0.9 \% \\
3 & 3.5 \% & 0.9 \% \\
10 & 39.5 \% & 1.0 \% \\
\end{longtable}
}

\subsubsection{Table S6d. Driver-set connectivity}\label{table-s6d.-driver-set-connectivity}

{\def\LTcaptype{none} % do not increment counter
\begin{longtable}[]{@{}
  >{\raggedright\arraybackslash}p{(\linewidth - 4\tabcolsep) * \real{0.3333}}
  >{\raggedright\arraybackslash}p{(\linewidth - 4\tabcolsep) * \real{0.3333}}
  >{\raggedright\arraybackslash}p{(\linewidth - 4\tabcolsep) * \real{0.3333}}@{}}
\toprule\noalign{}
\begin{minipage}[b]{\linewidth}\raggedright
set (same size)
\end{minipage} & \begin{minipage}[b]{\linewidth}\raggedright
in largest SCC
\end{minipage} & \begin{minipage}[b]{\linewidth}\raggedright
excitatory edges per node
\end{minipage} \\
\midrule\noalign{}
\endhead
\bottomrule\noalign{}
\endlastfoot
sustain drivers & 96 \% & 31 \\
top-degree hubs & 40 \% & 9.4 \\
random & 0.4 \% & 0.1 \\
\end{longtable}
}

\subsubsection{Table S6e. Targeted 20-s block on held-out seeds}\label{table-s6e.-targeted-20-s-block-on-held-out-seeds}

{\def\LTcaptype{none} % do not increment counter
\begin{longtable}[]{@{}
  >{\raggedright\arraybackslash}p{(\linewidth - 10\tabcolsep) * \real{0.1667}}
  >{\raggedright\arraybackslash}p{(\linewidth - 10\tabcolsep) * \real{0.1667}}
  >{\raggedright\arraybackslash}p{(\linewidth - 10\tabcolsep) * \real{0.1667}}
  >{\raggedright\arraybackslash}p{(\linewidth - 10\tabcolsep) * \real{0.1667}}
  >{\raggedright\arraybackslash}p{(\linewidth - 10\tabcolsep) * \real{0.1667}}
  >{\raggedright\arraybackslash}p{(\linewidth - 10\tabcolsep) * \real{0.1667}}@{}}
\toprule\noalign{}
\begin{minipage}[b]{\linewidth}\raggedright
anoxia
\end{minipage} & \begin{minipage}[b]{\linewidth}\raggedright
sustain drivers
\end{minipage} & \begin{minipage}[b]{\linewidth}\raggedright
earliest igniters
\end{minipage} & \begin{minipage}[b]{\linewidth}\raggedright
top-degree hubs
\end{minipage} & \begin{minipage}[b]{\linewidth}\raggedright
random
\end{minipage} & \begin{minipage}[b]{\linewidth}\raggedright
whole brain
\end{minipage} \\
\midrule\noalign{}
\endhead
\bottomrule\noalign{}
\endlastfoot
30 s: smallest rescuing set & \textbf{0.5 \% (696 neurons), 5/5 seeds, FII\_norm 0.947 {[}0.943, 0.951{]}} & \textbf{0.5 \%, 5/5, 0.945} & 2 \% (2,785), 5/5 & none up to 20 \% (best 0.33) & 1.000 \\
180 s: smallest rescuing set & \textbf{20 \% (27,851), 5/5, 0.979} (10 \%: 0/5) & \textbf{20 \%, 5/5} (10 \%: 0/5) & none up to 20 \% (0/5) & none up to 20 \% (0/5) & 0.979 \\
\end{longtable}
}

\subsubsection{Table S6f. Dissection (rate in last 10 s of a 40-s manipulation, Hz/neuron, 30 s / 180 s anoxia)}\label{table-s6f.-dissection-rate-in-last-10-s-of-a-40-s-manipulation-hzneuron-30-s-180-s-anoxia}

{\def\LTcaptype{none} % do not increment counter
\begin{longtable}[]{@{}
  >{\raggedright\arraybackslash}p{(\linewidth - 4\tabcolsep) * \real{0.3333}}
  >{\raggedright\arraybackslash}p{(\linewidth - 4\tabcolsep) * \real{0.3333}}
  >{\raggedright\arraybackslash}p{(\linewidth - 4\tabcolsep) * \real{0.3333}}@{}}
\toprule\noalign{}
\begin{minipage}[b]{\linewidth}\raggedright
manipulation
\end{minipage} & \begin{minipage}[b]{\linewidth}\raggedright
rate during manipulation
\end{minipage} & \begin{minipage}[b]{\linewidth}\raggedright
late rate after release
\end{minipage} \\
\midrule\noalign{}
\endhead
\bottomrule\noalign{}
\endlastfoot
all synapses off & 109 / 134 & 138 / 120 \\
excitatory synapses off & 14 / 120 & 38 / 91 \\
inhibitory synapses off & 133 / 135 & 111 / 115 \\
\textbf{K⁺ clamped at 5 mM} & \textbf{2.4 / 2.3} & 4.3 / 3.9 \\
\textbf{no spike-driven K⁺ release (k\_spk = 0)} & \textbf{2.4 / 2.3} & 4.3 / 4.1 \\
no SD K⁺ release (k\_rel = 0) & 56 / 153 & 129 / 25 \\
\end{longtable}
}

\subsubsection{Table S6g. Basin boundary: minimal K⁺ kick reaching the high-rate state vs tonic offset δ}\label{table-s6g.-basin-boundary-minimal-k-kick-reaching-the-high-rate-state-vs-tonic-offset-ux3b4}

{\def\LTcaptype{none} % do not increment counter
\begin{longtable}[]{@{}
  >{\raggedright\arraybackslash}p{(\linewidth - 14\tabcolsep) * \real{0.1250}}
  >{\raggedright\arraybackslash}p{(\linewidth - 14\tabcolsep) * \real{0.1250}}
  >{\raggedright\arraybackslash}p{(\linewidth - 14\tabcolsep) * \real{0.1250}}
  >{\raggedright\arraybackslash}p{(\linewidth - 14\tabcolsep) * \real{0.1250}}
  >{\raggedright\arraybackslash}p{(\linewidth - 14\tabcolsep) * \real{0.1250}}
  >{\raggedright\arraybackslash}p{(\linewidth - 14\tabcolsep) * \real{0.1250}}
  >{\raggedright\arraybackslash}p{(\linewidth - 14\tabcolsep) * \real{0.1250}}
  >{\raggedright\arraybackslash}p{(\linewidth - 14\tabcolsep) * \real{0.1250}}@{}}
\toprule\noalign{}
\begin{minipage}[b]{\linewidth}\raggedright
offset δ
\end{minipage} & \begin{minipage}[b]{\linewidth}\raggedright
+6 mV
\end{minipage} & \begin{minipage}[b]{\linewidth}\raggedright
+4 mV
\end{minipage} & \begin{minipage}[b]{\linewidth}\raggedright
+2 mV
\end{minipage} & \begin{minipage}[b]{\linewidth}\raggedright
0 mV
\end{minipage} & \begin{minipage}[b]{\linewidth}\raggedright
−2 mV
\end{minipage} & \begin{minipage}[b]{\linewidth}\raggedright
−4 mV
\end{minipage} & \begin{minipage}[b]{\linewidth}\raggedright
−6 mV
\end{minipage} \\
\midrule\noalign{}
\endhead
\bottomrule\noalign{}
\endlastfoot
minimal kick & 6 mM & 7 mM & 8 mM & 9 mM & 12 mM & 15 mM & none up to 30 mM \\
analytic K\_th & 5.3 & 6.1 & 6.9 & 7.9 & 9.0 & 10.2 & 11.7 mM \\
\end{longtable}
}

\begin{figure}
\centering
\includegraphics[width=1\linewidth,height=\textheight,keepaspectratio,alt={Figure S3. (a) Dead fraction after 180-s anoxia vs excitotoxic gain k\_exc (n = 5 seeds per point). (b) Paired Δ dead, graded − abrupt, no block. (c) Paired Δ FII\_norm, graded + block − abrupt + matched block. Dots = means with paired-bootstrap 95\% CIs.}]{figures/FigS3_excitotoxic_H3.png}
\caption{\textbf{Figure S3.} (a) Dead fraction after 180-s anoxia vs excitotoxic gain k\_exc (n = 5 seeds per point). (b) Paired Δ dead, graded − abrupt, no block. (c) Paired Δ FII\_norm, graded + block − abrupt + matched block. Dots = means with paired-bootstrap 95\% CIs.}
\end{figure}

\subsubsection{S6h. Transmitter-sign robustness (antennal-lobe local neurons inhibitory)}\label{s6h.-transmitter-sign-robustness-antennal-lobe-local-neurons-inhibitory}

Variant A sets every neuron with class ALLN or primary type matching lLN/vLN/v2LN/ilLN/il2LN to inhibitory. In FAFB the mask covers 431 neurons; 284 were already inhibitory and 147 changed sign (90 predicted cholinergic, 57 predicted monoaminergic). The same rule changed 147 neurons in the FAFB-proxy model and 219 in MCNS.

\textbf{Table S6h-1. FAFB full model, variant A, 30-s anoxia, naive reperfusion (per seed; local runs).} FII is raw (vs the default-sign intact baseline; no variant-A baseline was run); FII\_norm ≈ FII/0.9940. Runaway rule: late rate \textgreater{} 0.5 Hz/neuron or FSA \textgreater{} 0.05. Default-sign comparison (dynamics arm C, 5 seeds): FII\_norm −0.058, runaway 5/5; anatomy runs N\_anat\_D30 s0--s4: mean observation rate 91.3--104.3 Hz/neuron, 64.0--69.5 \% of neurons ≥ 20 Hz.

{\def\LTcaptype{none} % do not increment counter
\begin{longtable}[]{@{}
  >{\raggedright\arraybackslash}p{(\linewidth - 18\tabcolsep) * \real{0.1000}}
  >{\raggedright\arraybackslash}p{(\linewidth - 18\tabcolsep) * \real{0.1000}}
  >{\raggedright\arraybackslash}p{(\linewidth - 18\tabcolsep) * \real{0.1000}}
  >{\raggedright\arraybackslash}p{(\linewidth - 18\tabcolsep) * \real{0.1000}}
  >{\raggedright\arraybackslash}p{(\linewidth - 18\tabcolsep) * \real{0.1000}}
  >{\raggedright\arraybackslash}p{(\linewidth - 18\tabcolsep) * \real{0.1000}}
  >{\raggedright\arraybackslash}p{(\linewidth - 18\tabcolsep) * \real{0.1000}}
  >{\raggedright\arraybackslash}p{(\linewidth - 18\tabcolsep) * \real{0.1000}}
  >{\raggedright\arraybackslash}p{(\linewidth - 18\tabcolsep) * \real{0.1000}}
  >{\raggedright\arraybackslash}p{(\linewidth - 18\tabcolsep) * \real{0.1000}}@{}}
\toprule\noalign{}
\begin{minipage}[b]{\linewidth}\raggedright
run
\end{minipage} & \begin{minipage}[b]{\linewidth}\raggedright
seed
\end{minipage} & \begin{minipage}[b]{\linewidth}\raggedright
dead
\end{minipage} & \begin{minipage}[b]{\linewidth}\raggedright
late rate (Hz/neuron)
\end{minipage} & \begin{minipage}[b]{\linewidth}\raggedright
FSA
\end{minipage} & \begin{minipage}[b]{\linewidth}\raggedright
FII (raw)
\end{minipage} & \begin{minipage}[b]{\linewidth}\raggedright
FII/0.9940
\end{minipage} & \begin{minipage}[b]{\linewidth}\raggedright
surge peak (Hz)
\end{minipage} & \begin{minipage}[b]{\linewidth}\raggedright
K⁺ peak (mM)
\end{minipage} & \begin{minipage}[b]{\linewidth}\raggedright
runaway
\end{minipage} \\
\midrule\noalign{}
\endhead
\bottomrule\noalign{}
\endlastfoot
S\_A\_D30\_naive\_s0 & 0 & 0.000 & 0.168 & 0.0050 & 0.915 & 0.920 & 68.4 & 44.9 & no \\
S\_A\_D30\_naive\_s1 & 1 & 0.000 & 0.160 & 0.0049 & 0.928 & 0.933 & 69.6 & 45.0 & no \\
S\_A\_D30\_naive\_s2 & 2 & 0.000 & 0.035 & 0.0008 & 0.936 & 0.942 & 68.9 & 45.0 & no \\
\end{longtable}
}

{\def\LTcaptype{none} % do not increment counter
\begin{longtable}[]{@{}
  >{\raggedright\arraybackslash}p{(\linewidth - 8\tabcolsep) * \real{0.2000}}
  >{\raggedright\arraybackslash}p{(\linewidth - 8\tabcolsep) * \real{0.2000}}
  >{\raggedright\arraybackslash}p{(\linewidth - 8\tabcolsep) * \real{0.2000}}
  >{\raggedright\arraybackslash}p{(\linewidth - 8\tabcolsep) * \real{0.2000}}
  >{\raggedright\arraybackslash}p{(\linewidth - 8\tabcolsep) * \real{0.2000}}@{}}
\toprule\noalign{}
\begin{minipage}[b]{\linewidth}\raggedright
anatomy run
\end{minipage} & \begin{minipage}[b]{\linewidth}\raggedright
seed
\end{minipage} & \begin{minipage}[b]{\linewidth}\raggedright
mean observation rate (Hz/neuron)
\end{minipage} & \begin{minipage}[b]{\linewidth}\raggedright
fraction of neurons ≥ 20 Hz
\end{minipage} & \begin{minipage}[b]{\linewidth}\raggedright
median first spike (ms)
\end{minipage} \\
\midrule\noalign{}
\endhead
\bottomrule\noalign{}
\endlastfoot
SN\_A\_anat\_D30\_s0 & 0 & 0.172 & 0.0028 & 105 \\
SN\_A\_anat\_D30\_s1 & 1 & 0.160 & 0.0027 & 105 \\
SN\_A\_anat\_D30\_s2 & 2 & 0.035 & 0.0005 & 105 \\
\end{longtable}
}

\textbf{Table S6h-2. MCNS and FAFB-proxy, default signs vs antennal-lobe local neurons inhibitory (\texttt{\_alln}), 5 seeds per condition.}

{\def\LTcaptype{none} % do not increment counter
\begin{longtable}[]{@{}
  >{\raggedright\arraybackslash}p{(\linewidth - 14\tabcolsep) * \real{0.1250}}
  >{\raggedright\arraybackslash}p{(\linewidth - 14\tabcolsep) * \real{0.1250}}
  >{\raggedright\arraybackslash}p{(\linewidth - 14\tabcolsep) * \real{0.1250}}
  >{\raggedright\arraybackslash}p{(\linewidth - 14\tabcolsep) * \real{0.1250}}
  >{\raggedright\arraybackslash}p{(\linewidth - 14\tabcolsep) * \real{0.1250}}
  >{\raggedright\arraybackslash}p{(\linewidth - 14\tabcolsep) * \real{0.1250}}
  >{\raggedright\arraybackslash}p{(\linewidth - 14\tabcolsep) * \real{0.1250}}
  >{\raggedright\arraybackslash}p{(\linewidth - 14\tabcolsep) * \real{0.1250}}@{}}
\toprule\noalign{}
\begin{minipage}[b]{\linewidth}\raggedright
model
\end{minipage} & \begin{minipage}[b]{\linewidth}\raggedright
condition
\end{minipage} & \begin{minipage}[b]{\linewidth}\raggedright
n
\end{minipage} & \begin{minipage}[b]{\linewidth}\raggedright
outcomes
\end{minipage} & \begin{minipage}[b]{\linewidth}\raggedright
FII\_norm {[}95\% CI{]}
\end{minipage} & \begin{minipage}[b]{\linewidth}\raggedright
late rate (Hz/neuron)
\end{minipage} & \begin{minipage}[b]{\linewidth}\raggedright
FSA
\end{minipage} & \begin{minipage}[b]{\linewidth}\raggedright
dead
\end{minipage} \\
\midrule\noalign{}
\endhead
\bottomrule\noalign{}
\endlastfoot
FAFBproxy & D180\_T22\_naive & 5 & \{``runaway'': 5\} & 0.288 {[}0.086, 0.505{]} & 124.09 & 0.044 & 0.009 \\
FAFBproxy\_alln & D180\_T22\_naive & 5 & \{``runaway'': 5\} & -0.125 {[}-0.155, -0.093{]} & 126.62 & 0.743 & 0.009 \\
MCNS & D180\_T22\_naive & 5 & \{``runaway'': 5\} & 0.856 {[}0.810, 0.901{]} & 127.79 & 0.001 & 0.009 \\
MCNS\_alln & D180\_T22\_naive & 5 & \{``runaway'': 5\} & -0.135 {[}-0.234, -0.041{]} & 120.53 & 0.897 & 0.009 \\
FAFBproxy & D180\_T22\_reset & 5 & \{``recovered'': 5\} & 0.971 {[}0.961, 0.984{]} & 0.00 & 0.000 & 0.009 \\
FAFBproxy\_alln & D180\_T22\_reset & 5 & \{``recovered'': 5\} & 0.965 {[}0.948, 0.978{]} & 0.00 & 0.000 & 0.009 \\
MCNS & D180\_T22\_reset & 5 & \{``recovered'': 5\} & 0.997 {[}0.992, 1.001{]} & 0.00 & 0.000 & 0.009 \\
MCNS\_alln & D180\_T22\_reset & 5 & \{``recovered'': 5\} & 0.989 {[}0.970, 1.000{]} & 0.00 & 0.000 & 0.009 \\
FAFBproxy & D30\_T22\_naive & 5 & \{``runaway'': 5\} & 0.146 {[}0.057, 0.228{]} & 47.11 & 0.140 & 0.000 \\
FAFBproxy\_alln & D30\_T22\_naive & 5 & \{``runaway'': 5\} & 0.859 {[}0.852, 0.867{]} & 0.82 & 0.008 & 0.000 \\
MCNS & D30\_T22\_naive & 5 & \{``runaway'': 5\} & 0.050 {[}0.022, 0.078{]} & 90.86 & 0.084 & 0.000 \\
MCNS\_alln & D30\_T22\_naive & 5 & \{``runaway'': 5\} & 1.000 {[}0.996, 1.002{]} & 3.00 & 0.022 & 0.000 \\
FAFBproxy & D30\_T22\_reset & 5 & \{``recovered'': 5\} & 1.002 {[}1.001, 1.003{]} & 0.00 & 0.000 & 0.000 \\
FAFBproxy\_alln & D30\_T22\_reset & 5 & \{``recovered'': 5\} & 1.002 {[}1.001, 1.003{]} & 0.00 & 0.000 & 0.000 \\
MCNS & D30\_T22\_reset & 5 & \{``recovered'': 5\} & 0.999 {[}0.994, 1.001{]} & 0.00 & 0.000 & 0.000 \\
MCNS\_alln & D30\_T22\_reset & 5 & \{``recovered'': 5\} & 1.000 {[}0.997, 1.002{]} & 0.00 & 0.000 & 0.000 \\
FAFBproxy & base & 5 & \{``recovered'': 5\} & 1.003 {[}1.001, 1.003{]} & & 0.000 & 0.000 \\
FAFBproxy\_alln & base & 5 & \{``recovered'': 5\} & 1.003 {[}1.001, 1.003{]} & & 0.000 & 0.000 \\
MCNS & base & 10 & \{``recovered'': 10\} & 0.999 {[}0.996, 1.001{]} & & 0.000 & 0.000 \\
MCNS\_alln & base & 5 & \{``recovered'': 5\} & 1.000 {[}0.997, 1.002{]} & & 0.000 & 0.000 \\
\end{longtable}
}

{\def\LTcaptype{none} % do not increment counter
\begin{longtable}[]{@{}
  >{\raggedright\arraybackslash}p{(\linewidth - 2\tabcolsep) * \real{0.5000}}
  >{\raggedright\arraybackslash}p{(\linewidth - 2\tabcolsep) * \real{0.5000}}@{}}
\toprule\noalign{}
\begin{minipage}[b]{\linewidth}\raggedright
contrast (AL local neurons inhibitory − default signs; closed-loop − 20-s block)
\end{minipage} & \begin{minipage}[b]{\linewidth}\raggedright
Δ FII\_norm {[}95\% CI{]}
\end{minipage} \\
\midrule\noalign{}
\endhead
\bottomrule\noalign{}
\endlastfoot
MCNS\_alln\_minus\_main\_D30\_naive & +0.949 {[}+0.918, +0.977{]} \\
FAFBproxy\_alln\_minus\_main\_D30\_naive & +0.714 {[}+0.632, +0.801{]} \\
MCNS\_alln\_minus\_main\_D180\_naive & -0.991 {[}-1.119, -0.883{]} \\
FAFBproxy\_alln\_minus\_main\_D180\_naive & -0.413 {[}-0.622, -0.192{]} \\
MCNS\_closedloop\_minus\_reset20\_D180 & -0.071 {[}-0.081, -0.058{]} \\
FAFBproxy\_closedloop\_minus\_reset20\_D180 & -0.674 {[}-0.742, -0.624{]} \\
\end{longtable}
}

Variants B (all low-confidence calls inhibitory) and C (neurons without a call silenced) were prepared in code but not run. FAFB variant A was not run at 180 s, with the reset, for bistability kicks, or for targeted blocks.

\subsection{S7. Cross-connectome dynamic replication (MCNS v1.0; BANC v888 not analysed)}\label{s7.-cross-connectome-dynamic-replication-mcns-v1.0-banc-v888-not-analysed}

Source: \texttt{research/followup/replication/REPLICATION\_RESULTS.md} (211 analysed runs: MCNS 146, FAFB-proxy 65). MCNS: 166,700 neurons, 6,242,085 edges (≥5 synapses), w\_syn 0.0538 mV (same distance below the reverberation threshold as FAFB). Proxy substitutions: demand = total synapse count; K⁺ coupling = connection-based Jaccard 9-nearest-neighbour neuropil graph. A dense shared-neuron K⁺ adjacency was numerically unstable (spurious oscillation and runaway in sham runs) and was replaced; the FAFB dynamics track used the stable geometric adjacency and is unaffected. Responsive probes only (sugar → MN9 is silent in MCNS); noise ceilings MCNS 0.9972, FAFB-proxy 0.9922.

\subsubsection{Table S7a. Condition summary}\label{table-s7a.-condition-summary}

{\def\LTcaptype{none} % do not increment counter
\begin{longtable}[]{@{}
  >{\raggedright\arraybackslash}p{(\linewidth - 6\tabcolsep) * \real{0.2500}}
  >{\raggedright\arraybackslash}p{(\linewidth - 6\tabcolsep) * \real{0.2500}}
  >{\raggedright\arraybackslash}p{(\linewidth - 6\tabcolsep) * \real{0.2500}}
  >{\raggedright\arraybackslash}p{(\linewidth - 6\tabcolsep) * \real{0.2500}}@{}}
\toprule\noalign{}
\begin{minipage}[b]{\linewidth}\raggedright
condition
\end{minipage} & \begin{minipage}[b]{\linewidth}\raggedright
\textbf{MCNS} FII\_norm {[}95 \% CI{]} · dead · outcomes
\end{minipage} & \begin{minipage}[b]{\linewidth}\raggedright
FAFB-proxy
\end{minipage} & \begin{minipage}[b]{\linewidth}\raggedright
FAFB (dynamics track)
\end{minipage} \\
\midrule\noalign{}
\endhead
\bottomrule\noalign{}
\endlastfoot
30 s anoxia, naive & 0.050 {[}0.022, 0.078{]} · dead 0.000 · runaway 5/5 & 0.146 {[}0.057, 0.228{]} · dead 0.000 · runaway 5/5 & -0.058 {[}-0.180, 0.153{]} · runaway 5/5 \\
30 s, 20-s block & 0.999 {[}0.994, 1.001{]} · dead 0.000 · recovered 5/5 & 1.002 {[}1.001, 1.003{]} · dead 0.000 · recovered 5/5 & 1.001 {[}1.000, 1.001{]} · recovered 5/5 \\
180 s, naive & 0.856 {[}0.810, 0.901{]} · dead 0.009 · runaway 5/5 & 0.288 {[}0.086, 0.505{]} · dead 0.009 · runaway 5/5 & 0.084 {[}-0.118, 0.410{]} · runaway 5/5 \\
180 s, 20-s block & 0.997 {[}0.992, 1.001{]} · dead 0.009 · recovered 5/5 & 0.971 {[}0.961, 0.984{]} · dead 0.009 · recovered 5/5 & 0.953 {[}0.936, 0.967{]} · recovered 5/5 \\
180 s @17 °C, naive & 0.915 {[}0.874, 0.943{]} · dead 0.001 · runaway 5/5 & -- & -- \\
180 s @17 °C, block & 0.992 {[}0.984, 0.998{]} · dead 0.001 · recovered 5/5 & -- & -- \\
180 s @29 °C, naive & 0.483 {[}0.374, 0.592{]} · dead 0.128 · runaway 5/5 & -- & -- \\
180 s @29 °C, block & 0.806 {[}0.738, 0.868{]} · dead 0.127 · recovered 3/5, structural\_loss 2/5 & -- & -- \\
360 s, naive & -0.033 {[}-0.275, 0.255{]} · dead 0.251 · runaway 5/5 & -- & 0.085 {[}-0.120, 0.311{]} · runaway 5/5 \\
360 s, block & 0.759 {[}0.691, 0.833{]} · dead 0.246 · structural\_loss 5/5 & 0.613 {[}0.562, 0.671{]} · dead 0.246 · structural\_loss 5/5 & 0.578 {[}0.559, 0.602{]} · structural\_loss 5/5 \\
360 s hub 1 \%, naive & 0.235 {[}0.042, 0.396{]} · dead 0.249 · runaway 5/5 & -- & -0.218 {[}-0.286, -0.093{]} · runaway 3/3 \\
360 s hub 1 \%, block & 0.960 {[}0.948, 0.973{]} · dead 0.244 · structural\_loss 5/5 & 0.707 {[}0.678, 0.737{]} · dead 0.244 · structural\_loss 5/5 & 0.703 {[}0.668, 0.743{]} · structural\_loss 3/3 \\
360 s hub 10 \%, naive & -0.137 {[}-0.197, -0.071{]} · dead 0.223 · runaway 5/5 & -- & 0.342 {[}0.300, 0.378{]} · runaway 3/3 \\
360 s hub 10 \%, block & 0.977 {[}0.972, 0.980{]} · dead 0.219 · recovered 5/5 & -- & 0.906 {[}0.845, 0.944{]} · recovered 3/3 \\
360 s random 1 \%, block & 0.769 {[}0.704, 0.837{]} · dead 0.244 · structural\_loss 5/5 & -- & 0.568 {[}0.530, 0.591{]} · structural\_loss 3/3 \\
360 s random 10 \%, block & 0.813 {[}0.775, 0.851{]} · dead 0.222 · structural\_loss 5/5 & -- & 0.588 {[}0.563, 0.637{]} · structural\_loss 3/3 \\
360 s whole-brain, block & 0.999 {[}0.994, 1.001{]} · dead 0.000 · recovered 5/5 & -- & 1.001 {[}1.001, 1.001{]} · recovered 3/3 \\
180 s, closed-loop release & 0.926 {[}0.919, 0.939{]} · dead 0.007 · runaway 5/5 & 0.296 {[}0.228, 0.346{]} · dead 0.007 · runaway 5/5 & -- \\
sham (0 s), naive & 1.001 {[}1.000, 1.001{]} · dead 0.000 · recovered 3/3 & -- & -- \\
sham (0 s), block & 1.001 {[}1.000, 1.001{]} · dead 0.000 · recovered 3/3 & -- & -- \\
\end{longtable}
}

\subsubsection{Table S7b. Hub-protection efficiency in MCNS (360-s anoxia + 20-s block)}\label{table-s7b.-hub-protection-efficiency-in-mcns-360-s-anoxia-20-s-block}

{\def\LTcaptype{none} % do not increment counter
\begin{longtable}[]{@{}
  >{\raggedright\arraybackslash}p{(\linewidth - 4\tabcolsep) * \real{0.3333}}
  >{\raggedright\arraybackslash}p{(\linewidth - 4\tabcolsep) * \real{0.3333}}
  >{\raggedright\arraybackslash}p{(\linewidth - 4\tabcolsep) * \real{0.3333}}@{}}
\toprule\noalign{}
\begin{minipage}[b]{\linewidth}\raggedright
protected set (MCNS, n neurons)
\end{minipage} & \begin{minipage}[b]{\linewidth}\raggedright
gain in FII\_norm vs none
\end{minipage} & \begin{minipage}[b]{\linewidth}\raggedright
gain per 1,000 protected
\end{minipage} \\
\midrule\noalign{}
\endhead
\bottomrule\noalign{}
\endlastfoot
top-degree 1 \% (1,667) & +0.201 {[}+0.127, +0.278{]} & 0.1208 \\
top-degree 10 \% (16,670) & +0.218 {[}+0.143, +0.293{]} & 0.0131 \\
random 1 \% & +0.010 {[}-0.089, +0.114{]} & 0.0061 \\
random 10 \% & +0.054 {[}-0.028, +0.140{]} & 0.0032 \\
whole CNS (166,700) & +0.239 {[}+0.165, +0.315{]} & 0.0014 \\
\end{longtable}
}

\subsubsection{Table S7c. Replication verdicts}\label{table-s7c.-replication-verdicts}

{\def\LTcaptype{none} % do not increment counter
\begin{longtable}[]{@{}
  >{\raggedright\arraybackslash}p{(\linewidth - 2\tabcolsep) * \real{0.5000}}
  >{\raggedright\arraybackslash}p{(\linewidth - 2\tabcolsep) * \real{0.5000}}@{}}
\toprule\noalign{}
\begin{minipage}[b]{\linewidth}\raggedright
test
\end{minipage} & \begin{minipage}[b]{\linewidth}\raggedright
verdict (MCNS; BANC not analysed)
\end{minipage} \\
\midrule\noalign{}
\endhead
\bottomrule\noalign{}
\endlastfoot
(a) runaway attractor after naive reperfusion & \textbf{replicates} (5/5 at 30, 180, 360 s and 17/22/29 °C; continuous at 30 s, episodic at ≥180 s) \\
(b) 20-s block reset rescue & \textbf{replicates} (5/5 recovered at 30 and 180 s) \\
(c) hub-protection efficiency & \textbf{replicates, stronger} (\textasciitilde84× per neuron vs FAFB \textasciitilde30×; random controls null) \\
(d) temperature effect & \textbf{replicates in direction} (17 °C − 29 °C = +0.185); magnitude not comparable to FAFB at matched D \\
VNC vs brain & the VNC crosses 1 Hz/neuron first and is the most persistent part of the attractor \\
\end{longtable}
}

\textbf{BANC.} A BANC sweep (146 result files on the Modal Volume \texttt{flysim-data:/results\_rep\_v4/BANC*}) completed but was not retrieved or analysed because further cloud work was stopped (job recorded as failed; cost unknown, ≤ USD 1.76 by the container cap). MANC was not simulated dynamically.

\subsection{S8. Recoverability index (cross-validated; FAFB)}\label{s8.-recoverability-index-cross-validated-fafb}

Source: \texttt{research/followup/index/tables/model\_results.json} and \texttt{univariate\_*.csv}; 1,118 of 1,209 dynamics runs processed (91 not processed); condition-grouped 5-fold cross-validation (4--20 repeats depending on the model) with cluster-bootstrap 95\% CIs. Indices were computed on default-sign graphs. Intact FAFB: R\_eff (mean-field) 19.62; confident-transmitter R\_eff 7.38; signed 17.05 (\texttt{index/data/fafb\_intact\_index.json}). \texttt{ecore100} = fraction of neurons in the excitatory bootstrap core (maximal set in which every member receives enough excitatory synapses from other members to sustain firing at an assumed 100 Hz); \texttt{recip\_n} = surviving reciprocal excitatory pairs relative to intact; \texttt{reach\_sim} = simulator Reach.

\subsubsection{Table S8a. Post-reset FII\_norm, non-runaway runs (top single indices)}\label{table-s8a.-post-reset-fii_norm-non-runaway-runs-top-single-indices}

{\def\LTcaptype{none} % do not increment counter
\begin{longtable}[]{@{}llll@{}}
\toprule\noalign{}
index & n & CV R² {[}95\% CI{]} & Spearman ρ \\
\midrule\noalign{}
\endhead
\bottomrule\noalign{}
\endlastfoot
ecore100 & 562 & 0.964 {[}0.940, 0.978{]} & 0.949 \\
ecore200 & 562 & 0.964 {[}0.937, 0.980{]} & 0.940 \\
recip\_n & 562 & 0.961 {[}0.934, 0.975{]} & 0.943 \\
reach\_sim & 562 & 0.938 {[}0.891, 0.965{]} & 0.857 \\
hub5 & 562 & 0.923 {[}0.856, 0.955{]} & 0.824 \\
hub10 & 562 & 0.922 {[}0.857, 0.957{]} & 0.843 \\
hub1 & 562 & 0.914 {[}0.851, 0.950{]} & 0.795 \\
gscc\_exc & 562 & 0.908 {[}0.807, 0.951{]} & 0.920 \\
\end{longtable}
}

\subsubsection{Table S8b. Runaway after naive reperfusion among viable networks (top single indices)}\label{table-s8b.-runaway-after-naive-reperfusion-among-viable-networks-top-single-indices}

{\def\LTcaptype{none} % do not increment counter
\begin{longtable}[]{@{}
  >{\raggedright\arraybackslash}p{(\linewidth - 8\tabcolsep) * \real{0.2000}}
  >{\raggedright\arraybackslash}p{(\linewidth - 8\tabcolsep) * \real{0.2000}}
  >{\raggedright\arraybackslash}p{(\linewidth - 8\tabcolsep) * \real{0.2000}}
  >{\raggedright\arraybackslash}p{(\linewidth - 8\tabcolsep) * \real{0.2000}}
  >{\raggedright\arraybackslash}p{(\linewidth - 8\tabcolsep) * \real{0.2000}}@{}}
\toprule\noalign{}
\begin{minipage}[b]{\linewidth}\raggedright
index
\end{minipage} & \begin{minipage}[b]{\linewidth}\raggedright
n (runaway)
\end{minipage} & \begin{minipage}[b]{\linewidth}\raggedright
a-priori AUC {[}95\% CI{]}
\end{minipage} & \begin{minipage}[b]{\linewidth}\raggedright
CV AUC
\end{minipage} & \begin{minipage}[b]{\linewidth}\raggedright
Spearman ρ
\end{minipage} \\
\midrule\noalign{}
\endhead
\bottomrule\noalign{}
\endlastfoot
exc\_in\_frac & 482 (457) & 0.532 {[}0.450, 0.653{]} & 0.529 & 0.025 \\
rich1\_n & 482 (457) & 0.526 {[}0.484, 0.570{]} & 0.525 & 0.021 \\
EI\_syn & 482 (457) & 0.510 {[}0.466, 0.549{]} & 0.488 & 0.008 \\
ipr\_exc & 482 (457) & 0.479 {[}0.310, 0.790{]} & 0.297 & 0.016 \\
loop3\_super & 482 (457) & 0.364 {[}0.174, 0.738{]} & 0.297 & -0.109 \\
loop2\_super & 482 (457) & 0.379 {[}0.194, 0.742{]} & 0.263 & -0.098 \\
reach\_sim & 482 (457) & 0.351 {[}0.158, 0.800{]} & 0.254 & -0.118 \\
Reff\_mf\_known & 482 (457) & 0.342 {[}0.146, 0.726{]} & 0.254 & -0.124 \\
\end{longtable}
}

Combined models: post-reset FII random forest CV R² 0.986 {[}0.974, 0.992{]}; recovered-after-reset logistic CV AUC 0.999; runaway among viable networks random forest CV AUC 0.61 {[}0.38, 0.76{]} (chance). Leave-arm-out transfer of \texttt{ecore100} for post-reset FII: arm D R² 0.89, arm H 0.84, arm E (ageing) −0.50. The planned BANC/MCNS/MANC recoverability maps were started (FAFB random-removal grid to f = 0.6 only) and not completed; no standalone INDEX\_RESULTS report was written.

\subsection{S9. Deviations from the preregistration and task plan (all tracks)}\label{s9.-deviations-from-the-preregistration-and-task-plan-all-tracks}

\begin{enumerate}
\def\labelenumi{\arabic{enumi}.}
\tightlist
\item
  \textbf{Seeds and stimulus rates.} 3--5 seeds per condition (preregistered ≥10); one rate per probe (preregistered 3). Baseline used 10 seeds.
\item
  \textbf{Spontaneous-activity metric.} The Shiu model is silent at rest, so FSA\_baseline = 0 and the preregistered ratio is undefined. H5 used an evoked-neuron count (Act).
\item
  \textbf{Brain-death definition as run.} Spontaneous fraction \textless{} 0.01 and no readout response \textgreater{} 8 Hz (preregistered: \textless{} 2 SD of noise).
\item
  \textbf{Calibration choices made after inspecting runs.} J50 60 → 600 s, τ\_ATP 20 → 30 s, q10\_injury 2.5 added, w\_syn 0.275 → 0.185. All are reported with the rejected alternatives.
\item
  \textbf{H1 dynamical test} was fitted on FII vs dead fraction from insult runs, not on removal curves.
\item
  \textbf{H2(a)} (death-time correlation with demand) was not computed.
\item
  \textbf{H3 protocols} were not matched for integrated supply.
\item
  \textbf{Temperatures} 20 and 26 °C were not run (17/22/29/33 °C were).
\item
  \textbf{Reach in dynamics} used the ≥5 matrix only, with no cut-off sensitivity. Ageing synapse loss is not reflected in Reach.
\item
  \textbf{LT50} was extrapolated, not simulated (≤1 h anoxia simulated).
\item
  \textbf{Structural removal orders} were below the preregistered ≥20 for several regimes (FAFB random 13, metabolic\_area 10, metabolic\_syn 5, degree 3 + 2 fine-grid, focal 5, matched 3; BANC 15/8/3/4/2; MCNS 6/6/2/3/2; MANC full) because of CPU contention. The f grid was reduced from 24 to 13 points mid-run, L∞ was capped at 30 hops for later orders, and GSCC was computed for orders 0--2 only. The primary metric samples 128 source neurons.
\item
  \textbf{Structural controls reduced.} One rewired null (FAFB, MANC only); ≥1 cut-off only for MANC (FAFB/MCNS files are floored at 5).
\item
  \textbf{A CSV column-shift bug} (rows without GSCC) affected 143 FAFB, 78 BANC and 65 MCNS rows. These were repaired deterministically.
\item
  \textbf{MAOL v1.1} was not available locally and was not analysed (planned deviation).
\item
  \textbf{Mechanism follow-up of 605 runs not executed} (Modal spend limit): 40/60-s blocks for all LHS sets, the full k\_exc = 10 H3 grid, finer targeted-set brackets (0.2--0.3\% at 30 s; 12.5--17.5\% at 180 s), and a hysteresis ramp to −30 mV.
\item
  \textbf{Per-set FII baselines in arm L} used 4 seeds. Arm X k\_exc = 0 cells reused dynamics arm C runs (identical code path).
\item
  \textbf{Dynamic replication.} Earlier MCNS sweeps (v1, adjacency bug; v2, infrastructure failure) were discarded. The v4 sweep gave 211 analysed runs (MCNS + FAFB-proxy). The BANC v4 sweep (146 files) was not harvested or analysed. MCNS and FAFB-proxy use proxies for demand and K⁺ coupling. MANC was not simulated dynamically.
\item
  \textbf{Recoverability index.} 91 of 1,209 runs were not processed. Cross-connectome maps were not completed. Indices use default signs. No standalone report was written; results are taken from \texttt{model\_results.json}.
\item
  \textbf{Sign robustness.} Only the antennal-lobe local-neuron variant was run: FAFB full model 3 seeds at 30 s naive only (no variant-A intact baseline, so FII is relative to the default-sign baseline); MCNS and FAFB-proxy 5 seeds at 30/180 s naive and reset. Variants B and C were not run.
\item
  \textbf{Closed-loop release} was tested only in MCNS and the FAFB-proxy (180 s, default signs), not in the FAFB full model or across LHS sets.
\end{enumerate}

\subsection{S10. Negative and null results (reported in full)}\label{s10.-negative-and-null-results-reported-in-full}

\begin{enumerate}
\def\labelenumi{\arabic{enumi}.}
\tightlist
\item
  No sharp percolation threshold under random structural loss in any of four connectomes (H1 falsified; FAFB w 0.671).
\item
  No sharp threshold for FII vs dead fraction (ΔAIC +0.11 vs linear; width 1.56).
\item
  Graded, ramped or pulsed reperfusion was worse than abrupt reperfusion for identity and survival (H3 reversed). Adding an excitotoxic term did not change this for any comparison that included a reset.
\item
  Hyperpolarization (−10/−20 mV), GABA gain ×2--4, SD-release block, 3× clearance, noise, adaptation and cooling during reperfusion all failed to prevent the runaway.
\item
  Blocks of 2 and 5 s failed, and 10 s was unreliable at 360 s.
\item
  Preconditioning gave no protection (24.9\% vs 24.8\% dead).
\item
  Hyperthermia without hypoxia (33/37/40 °C) had no effect.
\item
  Metabolic support (reserve boost, clearance boost, mild hypometabolism) did not rescue the aged network.
\item
  ``Resuscitated but different'' did not exist outside the runaway state (0.0\% coverage with the reset).
\item
  No single cell type is a structural keystone (≤1.03\% reach loss).
\item
  Fly-like metabolic arrest alone did not produce hours-scale tolerance (extrapolated LT50 ≈ 1.51 h at J50 14,400 s vs fly 8.3 h).
\item
  No connectome index predicted naive-reperfusion runaway among viable networks (cross-validated best single-index AUC 0.53; combined 0.61 {[}0.38, 0.76{]}).
\item
  Targeted driver or hub blocks failed after 180 s of anoxia up to 10\% of the brain.
\item
  A closed-loop release at K⁺ \textless{} 6 mM and ATP \textgreater{} 0.9 (reached after 11 s) re-ignited the runaway in MCNS (runaway 5/5) and the FAFB-proxy.
\item
  With antennal-lobe local neurons inhibitory, the 30-s high-rate runaway did not form (FAFB 0/3 runaway; MCNS and FAFB-proxy kept a weak 3.0 / 0.82 Hz/neuron oscillation classed runaway by the preregistered rule). The 30-s driver-core result is therefore not established.
\end{enumerate}

\subsection{S11. Provenance and compute}\label{s11.-provenance-and-compute}

\begin{itemize}
\tightlist
\item
  Dynamics: Modal jobs a4ddd6e8-9ca3-4803-a889-fa5fdd58ccf9, 31c05682-cf05-4a51-b769-97b562f60172, 4d88396b-d4af-4b98-935c-ea62f206acfb, 937a47ef-8836-46d7-8764-a2a4d3655763 (smoke 79482286-adfb-4c8d-b57c-64d67726799e); 1,209 runs, 0 errors, 24.3 CPU-h.
\item
  Mechanism: smoke 900ff527-e669-4ffa-bcd0-5eed11e9297e; shard A 65780d67-c068-43a9-8cf0-1976a09783a9; shard B c6fe6ce5-53f5-4e4e-afe0-dc930dc25632; 2,061 runs, 0 errors. Follow-up 12a56f59-2ec3-4e3c-8934-253f564c9204 was killed by the spend limit and its results discarded.
\item
  Replication: calibration 6fd134b9-a8a3-4e02-8438-143947ce30ea; v1 MCNS 6e86553b-044e-4155-becd-21845a5eb524 (invalid); v2 04f516ec-6e0a-4889-86d3-d5d9a903a4f6 (infrastructure failure).
\item
  Replication (v4, analysed): smoke 14de3ae1-d0bf-412c-8841-c1a279754447 (MCNS 360-s hub arm, 45 runs); job A a9ef9894-0c76-4de2-89c9-6879f59167ff (166 runs). BANC job ef7d8554-1bbc-4a61-a027-11bb646177dc recorded as failed; 146 files on the Volume, not analysed.
\item
  Sign variant (FAFB full model): local runs, \texttt{research/followup/mechanism/raw/sign\_res/} (S\_A\_D30\_naive\_s0--2, SN\_A\_anat\_D30\_s0--2).
\item
  Recoverability index: local; \texttt{research/followup/index/}.
\item
  Structure: local 16-core host; pure NumPy bitset engine.
\item
  Identity check: the mechanism-track simulator with k\_exc = 0 reproduced dynamics run \texttt{C\_D30\_naive\_s0} exactly (FII 0.3617, late rate 122.998 Hz/neuron).
\item
  Run totals in the paper: 1,209 dynamics + 2,061 mechanism (784 of them in the LHS arm) + 211 replication + 6 FAFB sign-variant = 3,487 analysed simulated condition--seed runs.
\end{itemize}

\end{document}
